\documentclass[12pt]{article}

\usepackage[utf8]{inputenc}
\usepackage[T1]{fontenc}
\usepackage{lmodern}
\usepackage{amsmath}
\usepackage{booktabs}
\usepackage{graphicx}
\usepackage{geometry}
\usepackage[hidelinks]{hyperref}
\usepackage{natbib}
\usepackage{longtable}

\geometry{margin=1in}
\graphicspath{{Paper/figures/}{figures/}}
\newcommand{\tablepath}{Paper/tables/}
\newcommand{\figurelang}{en}
\newif\ifpolicyreport
\policyreportfalse

\title{Beyond Formality: The Firm-Size Wage Premium in Colombia}
\author{
Oliver Pardo\thanks{Associate Professor, Departamento de Administración, Pontificia Universidad Javeriana; Director, Centro Javeriano de Competitividad. Corresponding author: \href{mailto:pardoo@javeriana.edu.co}{pardoo@javeriana.edu.co}.}
\and
Jorge Orozco\thanks{Research Consultant, Centro Javeriano de Competitividad, Pontificia Universidad Javeriana.}
}
\date{\today}

\begin{document}

\maketitle

\begin{abstract}
This paper documents the firm-size wage premium in Colombia using household survey data for 2008--2019 and 2021--2025. Our preferred specification interacts firm size with formality and controls for gender, age, age squared, education, occupational position, and sector-department-year fixed effects. Relative to solo workers of the same formality status, workers in firms with more than 100 employees earn about 37 percent more within informal employment and 17 percent more within formal employment. Both differences are statistically significant. The informal premium rises sharply out of solo work and the smallest units, but is not monotonic across all sizes. Within formal employment, workers in the largest category earn more than workers in every smaller category. The premium is consistent with labor-market frictions that may constrain the growth of small businesses and limit worker reallocation toward larger, potentially more productive firms.
\end{abstract}

\vspace{0.5em}
\noindent\textbf{Keywords:} Firm-size wage premium; informality; Colombia.

\vspace{0.5em}
\noindent\textbf{JEL Codes:} J31; J46; O17; L25.


\section{Introduction}

Firm size is a well-established correlate of wages in labor economics \citep{BrownMedoff1989,OiIdson1999,troske1999evidence}. In labor markets where many workers are own-account workers or employed in micro-firms, wage differences between small and large firms help identify where questions about firm growth, worker allocation, and wage inequality are most relevant. This paper asks whether Colombian workers in larger firms earn higher real hourly wages than observably similar workers in smaller firms, how much of that gap is explained by observable worker characteristics, whether the gap has increased or decreased over time, and whether the wage gap between small and large firms differs between formal and informal employment.

Our motivation to measure the firm-size wage premium in Colombia is both empirical and policy-oriented. A recent policy diagnosis argues that low-income households are disproportionately linked to own-account work, very small businesses, informality, and low-earning jobs, whereas higher-income households are more often connected to larger and more formal employers \citep{eslava2026crecimiento}. This diagnosis is suggestive, but it does not show whether firm size remains associated with higher wages after accounting for worker characteristics and formality. The observed link between larger employers and higher household income may instead reflect the fact that larger firms employ workers with different characteristics, or that larger firms are more likely to offer formal jobs. This distinction matters because the wage gap associated with firm size has not been systematically documented in Colombia. The question is especially relevant in Colombia because employment remains concentrated outside the largest-firm categories, while recent growth theory links the rise of large firms to economic development \citep{Chen2026}.

We study these questions using harmonized microdata from the Gran Encuesta Integrada de Hogares (GEIH) for 2008--2019 and 2021--2025. Our preferred specification interacts firm size with formality and controls for gender, age, age squared, education, occupational position, and sector-department-year fixed effects. It compares workers within the same formality group and accounts for differences between own-account work, wage employment, and other employment relationships. Simpler specifications document how the premium changes as controls are added and how it evolves over time. Because unobserved worker characteristics may still be correlated with firm size, the estimates are conditional wage gaps rather than causal effects of firm size. The literature offers competing explanations for a premium that survives observable controls, including unobserved sorting and employer wage-setting differences. Such a premium does not by itself establish allocative inefficiency \citep{BrownMedoff1989,OiIdson1999,AbowdKramarzMargolis1999,CardHeiningKline2013,ElBadaouiStroblWalsh2010,BalkanTumen2016}.

The evidence delivers five main results. First, employment is concentrated at the bottom of the firm-size distribution. Solo workers and firms with up to ten workers account for about two thirds of employment. Second, raw wage gaps are large across firm-size categories, sectors, education groups, and formality status, while the wage gap by sex is small by comparison. Third, in the preferred specification, informal workers in firms with 101 or more employees earn 37.2 percent more than comparable informal solo workers. The premium rises sharply across the smallest categories, but does not increase monotonically through the largest category. Fourth, formal workers in 101+ firms earn 16.5 percent more than comparable formal solo workers, a statistically significant difference. They also earn more than workers in every smaller formal category. The formal size premium is therefore smaller than the informal premium relative to solo work, but it is not absent. Fifth, the complementary dynamic exercise using the simpler common-premium specification shows a positive large-firm premium throughout the period. Its point estimates suggest an increase between 2008 and 2025, but the change is not statistically significant.

The rest of the paper is organized as follows. Section~2 describes the GEIH data, the harmonization of firm-size categories, and the construction of real hourly labor income. Section~3 presents the descriptive evidence on employment shares and wage gaps by firm size, formality, education, gender, sector, and year. Section~4 introduces the econometric specifications. Section~5 reports the preferred estimates, comparisons with simpler specifications, dynamic results, and the elasticity benchmark, contrasting each result with the relevant literature. Section~6 discusses interpretation, limitations, and policy implications. Section~7 concludes.

\section{Data}

The empirical analysis uses microdata from the Gran Encuesta Integrada de Hogares (GEIH), Colombia's nationally representative labor force survey produced by the National Administrative Department of Statistics (DANE). Our sample covers 2008--2019 and 2021--2025. We exclude 2020 because of pandemic-related data limitations \citep{DANE2020GEIH}. The data combine separate annual samples rather than following the same workers over time.

The sample is restricted to employed workers with positive labor income and valid weekly hours. Let \(i\) index workers and \(t\) survey years. Denote reported monthly labor income by \(MonthlyIncome_{it}\), weekly hours by \(WeeklyHours_{it}\), and the end-of-year Consumer Price Index by \(CPI_t\). Real hourly labor income in constant 2025 Colombian pesos, \(w_{it}\), is

\begin{equation}
  w_{it} =
  \frac{12 \times MonthlyIncome_{it}}{52 \times WeeklyHours_{it}}\times
  \frac{CPI_{2025}}{CPI_t}.    
\end{equation}

Firm size is measured using the worker-reported firm-size variable in GEIH. We recode responses into categories that are comparable across survey years: Solo, 2--3, 4--5, 6--10, 11--19, 20--30, 31--50, 51--100, and 101+. Solo identifies respondents who report working alone. It is distinct from occupational position, which identifies whether they are own-account workers, employees, or employers. For example, a household's sole domestic employee is classified as a solo worker. The newer questionnaire separates 101--200 and 201+ workers. We combine both categories into 101+ to maintain comparability with earlier years. The data do not contain employer identifiers, so we cannot link respondents who work for the same firm.

The education variable distinguishes no education, preschool, primary, lower secondary, upper secondary, and higher education, as well as a separate ``does not know/not reported'' response. The descriptive education comparisons combine no education and preschool into ``preschool or less,'' and the two secondary categories into ``secondary.'' Only these descriptive education comparisons omit the ``does not know/not reported'' group.

Formality is defined using pension contribution status. Workers who contribute to pensions are classified as formal, while workers who do not contribute are classified as informal. This definition follows the social-protection approach used in Latin American labor-market measurement, which identifies informality through lack of access to social-security benefits such as pension or health insurance \citep{WorldBankLACInformality2026}. This is a worker-level measure of formality, not a measure of the firm's legal registration.

We group economic-activity codes from successive versions of Colombia's industrial classification (CIIU) into a common set of broad sectors. The geographic variable identifies departments, treating Bogot\'a as a department-equivalent Capital District.
Occupational position distinguishes private-sector employees, government
employees, domestic workers, own-account workers, employers, unpaid family
workers, unpaid workers in other households, day laborers, and other workers.
We recode occupational position so that each code denotes the same employment relationship across survey questionnaires \citep{DANE2025GEIHQuestionnaire}.

The initial sample contains 5,353,440 worker observations pooled across survey years. Of these, 5,268,514 have complete information on the variables just described, together with sex and age. A recorded ``does not know/not reported'' education response is retained as its own category and is distinct from a missing value.


\section{Descriptive Statistics}

Table~\ref{tab:descriptive-core-size-2008-2025} establishes two raw facts behind the paper. First, mean real hourly income is substantially higher in large firms than in solo work. In 2008, solo workers earned 5,388 Colombian pesos at 2025 prices per hour on average, compared with 12,943 pesos among workers in firms with 101 or more workers, a ratio of 2.4 to 1. By 2025, the corresponding means were 6,534 and 16,859 pesos, a ratio of 2.6 to 1. Relative to workers in firms with 11--19 workers, 101+ workers earned 1.49 times as much in 2008 and 1.45 times as much in 2025. The all-worker mean increased from 7,752 to 10,295 pesos, equivalent to 1.7 percent annualized real growth between 2008 and 2025. Average weekly hours in the sample declined from 47.8 to 43.8 over the same period. This decline helps explain why real hourly income grew faster than implied monthly labor income per worker, which increased by 1.5 percent per year. Neither weekly hours nor implied monthly income is reported in the table.

Second, employment remains concentrated in solo work and the smallest firms. In 2025, solo workers accounted for 35 percent of employment. Together with workers in firms with 2--10 workers, they accounted for 61 percent. Firms with 11--100 workers jointly accounted for only 14 percent of employment. The 101+ category grew from 19 percent of employment in 2008 to 25 percent in 2025. This shift toward larger employers is consistent with the increasing prevalence of very large firms with economic development documented by \citet{Chen2026}.

\begin{table}[htbp]
  \centering
  \caption{Employment shares and mean real hourly income, 2008 vs. 2025}
  \label{tab:descriptive-core-size-2008-2025}
  \small
  \begin{tabular}{lrrrrrrr}
    \toprule
    & \multicolumn{2}{c}{Employment share} & \multicolumn{2}{c}{Mean hourly income} & \multicolumn{2}{c}{Wage ratio to solo} & \multicolumn{1}{c}{Wage growth} \\
    \cmidrule(lr){2-3} \cmidrule(lr){4-5} \cmidrule(lr){6-7} \cmidrule(lr){8-8}
    Firm size & 2008 & 2025 & 2008 & 2025 & 2008 & 2025 & Annualized \\
    \midrule
    Solo & 36.4\% & 34.6\% & 5,179 & 6,394 & 1.0 & 1.0 & 1.2\% \\
    2-3 & 19.1\% & 13.7\% & 5,871 & 7,097 & 1.1 & 1.1 & 1.1\% \\
    4-5 & 7.4\% & 7.0\% & 6,790 & 8,096 & 1.3 & 1.3 & 1.0\% \\
    6-10 & 6.0\% & 5.8\% & 7,481 & 9,426 & 1.4 & 1.5 & 1.4\% \\
    11-19 & 3.5\% & 3.9\% & 8,825 & 11,446 & 1.7 & 1.8 & 1.5\% \\
    20-30 & 3.3\% & 3.3\% & 9,175 & 10,990 & 1.8 & 1.7 & 1.1\% \\
    31-50 & 2.6\% & 3.0\% & 10,862 & 12,408 & 2.1 & 1.9 & 0.8\% \\
    51-100 & 3.1\% & 3.5\% & 9,918 & 13,039 & 1.9 & 2.0 & 1.6\% \\
    101+ & 18.5\% & 25.2\% & 12,757 & 16,549 & 2.5 & 2.6 & 1.5\% \\
    \addlinespace
    All workers & 100.0\% & 100.0\% & 7,525 & 10,103 & 1.5 & 1.6 & 1.7\% \\
    \bottomrule
  \end{tabular}
  \vspace{0.3em}
  \begin{minipage}{0.95\textwidth}
  \footnotesize
  Notes: Mean hourly income is expressed in constant 2025 pesos. Annualized income growth is compounded over 2008--2025.
  \end{minipage}
\end{table}


Table~\ref{tab:descriptive-wage-gaps-2025} puts the raw firm-size gap in
context by comparing it with remuneration gaps across sex, formality status,
education, and economic activities in 2025. The wage gap by sex is small in
these means: women's mean income is 1.01 times men's. Formal workers earn 2.43 times as much as informal workers,
close to the 2.6-to-1 ratio between workers
in 101+ firms and solo workers reported in
Table~\ref{tab:descriptive-core-size-2008-2025}.
Workers with higher education earn 4.08 times as much as workers with
preschool or less. Workers in information and communication earn 3.63 times
as much as workers in agriculture. The raw
firm-size gap is therefore far larger than the wage gap by sex, comparable to
the formality gap, and smaller than the largest education and sector gaps.

\begin{table}[!htbp]
  \centering
  \caption{Mean real hourly labor income gaps by worker group, 2025}
  \label{tab:descriptive-wage-gaps-2025}
  \scriptsize
  \begin{tabular}{lp{0.34\textwidth}rrrr}
    \toprule
    Dimension & Group & Employment share & Mean income & Gap & Ratio \\
    \midrule
    Sex & Men & 58.8\% & 10,027 & -0.8\% & 1.00 \\
     & Women & 41.2\% & 10,212 & 1.1\% & 1.02 \\
    \midrule
    Formality & Formal & 44.9\% & 14,952 & 48.0\% & 2.43 \\
     & Informal & 55.1\% & 6,150 & -39.1\% & 1.00 \\
    \midrule
    Education & Preschool or less & 2.9\% & 4,256 & -57.9\% & 1.00 \\
     & Primary & 18.1\% & 5,559 & -45.0\% & 1.31 \\
     & Secondary & 45.5\% & 7,010 & -30.6\% & 1.65 \\
     & Higher education & 33.5\% & 17,265 & 70.9\% & 4.06 \\
    \midrule
    Activity & A - Agriculture & 13.7\% & 5,870 & -41.9\% & 1.00 \\
     & T - Household employers & 3.1\% & 6,392 & -36.7\% & 1.09 \\
     & I - Accommodation and food & 7.6\% & 6,814 & -32.6\% & 1.16 \\
     & H - Transport & 7.8\% & 7,700 & -23.8\% & 1.31 \\
     & F - Construction & 6.9\% & 8,322 & -17.6\% & 1.42 \\
     & G - Commerce & 17.0\% & 8,423 & -16.6\% & 1.44 \\
     & R/S - Arts and other services & 5.4\% & 8,545 & -15.4\% & 1.46 \\
     & C - Manufacturing & 11.0\% & 9,495 & -6.0\% & 1.62 \\
     & D/E - Utilities and water & 1.3\% & 11,102 & 9.9\% & 1.89 \\
     & B - Mining & 1.3\% & 11,136 & 10.2\% & 1.90 \\
     & L/M/N - Business services & 9.0\% & 13,176 & 30.4\% & 2.24 \\
     & Q - Health & 4.4\% & 14,376 & 42.3\% & 2.45 \\
     & P - Education & 4.2\% & 19,335 & 91.4\% & 3.29 \\
     & K - Finance and insurance & 1.9\% & 21,221 & 110.0\% & 3.62 \\
     & O - Public administration & 3.6\% & 20,955 & 107.4\% & 3.57 \\
     & J - Information and communication & 1.7\% & 21,608 & 113.9\% & 3.68 \\
    \bottomrule
  \end{tabular}
  \vspace{0.3em}
  \begin{minipage}{0.95\textwidth}
  \footnotesize
  Notes: Employment shares are computed within each dimension. Mean income is real hourly labor income in constant 2025 pesos. The gap is the percent difference in each group's mean relative to the overall 2025 mean. The ratio divides each group's mean by the lowest mean within the same dimension. The sample excludes employed pension recipients. Sector rows omit extraterritorial organizations.
  \end{minipage}
\end{table}


Table~\ref{tab:descriptive-worker-characteristics-2025} shows why the raw firm-size wage premium cannot be attributed to firm size alone. In 2025, only 8.3 percent of solo workers were formal, compared with 46.6 percent in firms with 6--10 workers and 98.0 percent in firms with 101 or more workers. The share of workers with higher education also rises strongly with firm size, from 16.7 percent among solo workers to 62.9 percent in the 101+ category. The share of women varies less systematically across firm-size groups, while solo workers are older on average. These differences motivate the comparisons that follow, which ask whether the firm-size premium is visible within sex, formality, education, and sector groups.

\begin{table}[htbp]
  \centering
  \caption{Worker characteristics by firm size, 2025}
  \label{tab:descriptive-worker-characteristics-2025}
  \small
  \begin{tabular}{lrrrr}
    \toprule
    Firm size & Women & Formal & Higher education & Mean age \\
    \midrule
    Solo & 41.3\% & 8.3\% & 16.7\% & 44.5 \\
    2-3 & 34.6\% & 13.0\% & 17.0\% & 41.3 \\
    4-5 & 36.5\% & 24.2\% & 20.2\% & 37.6 \\
    6-10 & 40.5\% & 46.6\% & 29.9\% & 37.0 \\
    11-19 & 40.4\% & 64.4\% & 38.8\% & 37.2 \\
    20-30 & 42.7\% & 81.2\% & 45.4\% & 37.0 \\
    31-50 & 41.2\% & 90.4\% & 51.0\% & 37.5 \\
    51-100 & 43.7\% & 94.3\% & 53.4\% & 37.5 \\
    101+ & 45.8\% & 98.0\% & 62.9\% & 38.0 \\
    \bottomrule
  \end{tabular}
  \vspace{0.3em}
  \begin{minipage}{0.95\textwidth}
  \footnotesize
  Notes: Statistics use only 2025 observations and are weighted by GEIH expansion weights. Higher education corresponds to workers in the harmonized higher-education category.
  \end{minipage}
\end{table}


Figure~\ref{fig:income-formality-comparison} shows the raw relationship between firm size and mean real hourly income separately for formal and informal workers. The points report mean hourly income by firm-size category in 2008 and 2025, and the bubble size represents each category's employment share within the corresponding formality-year group. Formal workers earn more than informal workers and are much more concentrated in large firms. The wage pattern by firm size is still visible within informal workers: among them, mean hourly income rises from 5,677 pesos for solo workers to 12,813 pesos in firms with 101 or more workers in 2025. Among formal workers, the raw profile is closer to a U-shape: formal solo workers earned 15,990 pesos per hour in 2025, formal workers in firms with 4--5 workers earned 12,036 pesos, and formal workers in firms with 101 or more workers earned 16,944 pesos.

\begin{figure}[!htbp]
  \centering
  \includegraphics[width=0.95\textwidth]{fig69_\figurelang.png}
  \caption{Mean real hourly income by firm size and formality status, 2008 and 2025}
  \label{fig:income-formality-comparison}
\end{figure}

Table~\ref{tab:formal-worker-profile-2025} helps interpret the left panel of Figure~\ref{fig:income-formality-comparison}. Formal solo work is a heterogeneous category. By occupational position, 82.1 percent are own-account workers and 16.2 percent are domestic workers. By education, 50.9 percent have higher education. A further breakdown shows that, among formal solo workers with higher education, 95.5 percent are own-account workers, with legal and accounting services among the most common activities. Their mean hourly income is 21,826 pesos, compared with 9,950 pesos among formal solo workers without higher education. Formal solo work thus includes relatively well-paid independent professionals, but the category is not made up exclusively of them.

\begin{table}[htbp]
  \centering
  \caption{Profile of formal workers by firm-size group, 2025}
  \label{tab:formal-worker-profile-2025}
  \small
  \begin{tabular}{lrrr}
    \toprule
    Characteristic & Formal solo & Formal 2--100 & Formal 101+ \\
    \midrule
    \multicolumn{4}{l}{\textit{Employment shares}} \\
    Share of formal workers & 6.4\% & 38.7\% & 54.9\% \\
    \midrule
    \multicolumn{4}{l}{\textit{Remuneration}} \\
    Mean real hourly income (2025 pesos) & 15,678 & 12,443 & 16,633 \\
    \midrule
    \multicolumn{4}{l}{\textit{Personal characteristics}} \\
    Mean age & 45.0 & 38.3 & 37.9 \\
    Women & 43.3\% & 41.9\% & 45.7\% \\
    Higher education & 50.8\% & 48.2\% & 63.4\% \\
    \midrule
    \multicolumn{4}{l}{\textit{Status in employment}} \\
    Own-account workers & 82.1\% & 7.7\% & 9.6\% \\
    Domestic workers & 16.2\% & 0.9\% & 0.0\% \\
    \bottomrule
  \end{tabular}
  \vspace{0.3em}
  \begin{minipage}{0.95\textwidth}
  \footnotesize
  Notes: Employment shares use all formal workers as the denominator. Other percentages are within each firm-size group. The status-in-employment categories shown are not an exhaustive list. Higher education denotes education beyond secondary school.
  \end{minipage}
\end{table}


To compare wage distributions with the minimum wage, we first define an hourly benchmark for 2025. We divide the statutory monthly minimum of 1,423,500 pesos by 45 weekly hours and $52/12$ weeks per month. Figure~\ref{fig:minwage-distribution-formality-size} shows kernel densities of hourly income relative to this benchmark for each size-formality group. The horizontal axis uses a logarithmic scale, and the vertical line marks income equal to the minimum wage. Among formal workers, the categories from 4--5 to 51--100 workers are visibly concentrated around the minimum wage. Compared with these categories, formal solo workers have a broader distribution, and formal workers in 101+ firms have more mass at higher incomes. Informal workers are concentrated below the minimum wage, particularly in solo work and firms with 2--3 workers. The employment shares reported in the figure also show that 54.9 percent of formal workers are in 101+ firms, whereas 57.6 percent of informal workers are solo workers.

\begin{figure}[!htbp]
  \centering
  % Keep the density curves intact while replacing the legacy benchmark labels.
  \includegraphics[width=\textwidth,trim=0 24bp 0 45.6bp,clip]{fig77_\figurelang.png}
  \par\smallskip
  {\small\textbf{Hourly income / hourly minimum-wage benchmark (log scale)}\par}
  \caption{Smoothed hourly wage distributions relative to the minimum-wage benchmark by firm size and formality status, 2025}
  \label{fig:minwage-distribution-formality-size}
\end{figure}

Figure~\ref{fig:income-sex-comparison} shows that the raw firm-size pattern is
similar for men and women. In 2025, among men, solo workers earned 6,816 pesos
per hour, compared with 16,929 pesos in firms with 101 or more workers. Among
women, the corresponding means were 6,133 and 16,777 pesos.
Relative to each sex's solo-worker mean, the corresponding ratios are 2.5
and 2.7, respectively. The figure also shows that wage
differences between men and women within a given firm-size category are small
relative to the difference between solo work and large-firm employment.

\begin{figure}[!htbp]
  \centering
  \includegraphics[width=0.95\textwidth]{fig70_\figurelang.png}
  \caption{Mean real hourly income by firm size and sex, 2008 and 2025}
  \label{fig:income-sex-comparison}
\end{figure}

Figure~\ref{fig:income-education-comparison} shows that the raw firm-size
premium is visible within every education group in 2025. Among workers with preschool
or less, mean hourly income in 101+ firms is 2.1 times solo-worker income. The
corresponding 101+-to-solo ratio is 1.8 among workers with primary education,
1.5 among workers with secondary education, and 1.7 among workers with higher
education.

\begin{figure}[!htbp]
  \centering
  \includegraphics[width=\textwidth]{fig72_\figurelang.png}
  \caption{Mean real hourly income by firm size and education group, 2008 and 2025}
  \label{fig:income-education-comparison}
\end{figure}

\begin{figure}[p]
  \begin{minipage}{\textwidth}
Figure~\ref{fig:income-sector-comparison} shows that the raw firm-size wage
premium is widespread across economic activities in both 2008 and 2025.
Workers in firms with 101 or more workers earned more per hour than solo
workers in 14 of the 16 sectors in both years, and in all 16 sectors in 2025.
In agriculture, the ratio of mean hourly income in 101+ firms to solo-worker
income was approximately 1.9 in both years. In manufacturing, it rose from
2.35 to 2.63. The magnitude varies across sectors, and
wages do not rise at every size category.
  \end{minipage}
  \par\medskip
  \centering
  \includegraphics[width=\textwidth]{fig71_\figurelang.pdf}
  \caption{Mean real hourly income by firm size and sector, 2008 and 2025}
  \label{fig:income-sector-comparison}
\end{figure}

\FloatBarrier
The raw large-firm wage advantage is visible within each sex and education
group, in most sectors, and within formal and informal employment. Worker
composition also differs across firm sizes. The regression analysis examines
how much of these gaps remains after jointly controlling for worker
characteristics and occupational position, comparing workers within each
formality group and within the same sector, department, and year.

\section{Methodology}
\label{sec:methodology}

\subsection{Estimating the Firm-Size Wage Premium}
\label{subsec:preferred-specification}

%We begin with a regression of log real hourly labor income on firm-size indicators alone, with solo work as the reference category. We then add sector-year fixed effects, worker characteristics, and formality in sequence.

Let \(i\), \(s\), and \(t\) index workers, sectors, and survey years.
The variable \(w_{ist}\) denotes real hourly income.
Write \(D_{ig}\) for the indicator of size category \(g\), and \(g_0\) for
solo work. The vector \(X_i\) contains a female-worker indicator, age,
age squared, and education indicators. The formality indicator \(F_i\) equals one for pension
contributors and zero for noncontributors. With sector-year fixed effects
\(\lambda_{s\times t}\), we estimate
\begin{equation}\label{eq:common-premium}
  \ln(w_{ist}) = \alpha
  + \sum_{g\neq g_0} \beta_g D_{ig}
  + X_i'\gamma + \phi F_i
  + \lambda_{s\times t}+\varepsilon_{ist}.
\end{equation}
%All stages in this sequence use the same 5,268,514 complete observations and GEIH expansion weights. This sequence describes how the size coefficients change after accounting for observed differences between workers and sectors. Even after adding \(F\), the model constrains each size coefficient \(\beta_g\) to be the same for formal and informal workers.

%We next retain this common-premium restriction while adding geographic controls and occupational position. First, we add department-year fixed effects to Equation~(\ref{eq:common-premium}). We then replace the separate
%sector-year and department-year effects with sector-department-year fixed effects, before adding occupational-position fixed effects. 
Next, we add occupational-position fixed effects \(\rho_o\) and replace
the sector-year effects with sector-department-year fixed effects
\(\lambda_{s\times d\times t}\). The indices \(o\) and \(d\) denote
occupational positions and departments.
\begin{equation}\label{eq:common-premium-geography}
  \ln(w_{iosdt}) = \alpha
  + \sum_{g\neq g_0} \beta_g D_{ig}
  + X_i'\gamma + \phi F_i
  + \rho_o + \lambda_{s\times d\times t}
  + \varepsilon_{iosdt}.
\end{equation}
%The specifications with sector-department-year effects exclude 139 singleton observations, leaving 5,268,375 observations.

Our preferred specification allows each size coefficient to differ
between formal and informal workers by interacting the size indicators with
\(F_i\). For size \(g\), the log-wage gap relative to solo workers of the
same formality status is \(\beta_g\) for informal workers and
\(\beta_g+\theta_g\) for formal workers. We normalize
\(\beta_{g_0}=\theta_{g_0}=0\), making \(\phi\) the formal-informal
log-wage gap among solo workers.
\begin{equation}\label{eq:main}
  \ln(w_{iosdt}) = \alpha + \phi F_i
  + \sum_{g\neq g_0}(\beta_g+\theta_g F_i)D_{ig}
  + X_i'\gamma + \rho_o + \lambda_{s\times d\times t}
  + \varepsilon_{iosdt}.
\end{equation}
%The size coefficients compare firm-size categories within sector-department-year cells, separately by formality and conditional on the included worker controls and occupational fixed effects. Controlling for occupational position is relevant because solo work and employment in larger firms differ substantially in the prevalence of own-account and wage employment. %These controls describe conditional comparisons, rather than identifying the total effect of moving from own-account work to a salaried job.

%We estimate the preferred model with GEIH expansion weights using the same 5,268,375 observations as Equation~(\ref{eq:common-premium-geography}).

For an origin size \(g\) and a destination size \(h\), let
\(\Delta^F_{g\to h}\) and \(\Delta^I_{g\to h}\) denote the formal and
informal within-status log-wage differences. Let \(\Delta^{I\to F}_g\)
denote the formal-minus-informal log-wage difference at size \(g\). Then
\begin{align}
  \Delta^{F}_{g\to h} &= (\beta_h+\theta_h)-(\beta_g+\theta_g), \nonumber\\
  \Delta^{I}_{g\to h} &= \beta_h-\beta_g, \label{eq:preferred-contrasts}\\
  \Delta^{I\to F}_{g} &= \phi+\theta_g. \nonumber
\end{align}

For each estimated contrast \(\widehat{\Delta}\), we report the percentage gap
\(100[\exp(\widehat{\Delta})-1]\). This transformation compares
exponentiated conditional mean log incomes, not necessarily arithmetic mean
incomes as in the descriptive tables.

Standard errors are clustered by sector for Equation~(\ref{eq:common-premium})
and by sector-department for Equations~(\ref{eq:common-premium-geography})
and~(\ref{eq:main}). Tests of individual coefficients and contrasts are
two-sided. The appendix details the alternative clustering and wild cluster
bootstrap procedures.

\subsection{Changes over Time}

We extend Equation~(\ref{eq:main}) by allowing the firm-size premiums, the
formality coefficient, and all worker and occupational-position coefficients
to vary by year. Let \(\mathcal T\) denote the sample years and
\(\mathbf{1}\{t=\tau\}\) the indicator for year \(\tau\). The model is
\begin{equation}\label{eq:preferred-dynamics}
  \begin{aligned}
    \ln(w_{iosdt}) ={}&
    \sum_{\tau\in\mathcal T}\mathbf{1}\{t=\tau\}
    \left[
      \phi_\tau F_i
      + \sum_{g\neq g_0}(\beta_{g\tau}+\theta_{g\tau}F_i)D_{ig}
      + X_i'\gamma_\tau+\rho_{o\tau}
    \right] \\
    &+ \lambda_{s\times d\times t}+\varepsilon_{iosdt}.
  \end{aligned}
\end{equation}
We estimate annual premiums using separate regressions for each year.

To test the change between 2008 and 2025, we estimate both years jointly,
using the fully year-interacted specification in
Equation~(\ref{eq:preferred-dynamics}). Joint estimation retains the
covariance between the two years' estimates within sector-department clusters.
For informal workers, the null is
\(H_0:\beta_{g,2025}-\beta_{g,2008}=0\). For formal workers, it is
\(H_0:(\beta_{g,2025}+\theta_{g,2025})-(\beta_{g,2008}+\theta_{g,2008})=0\).
The tests are unadjusted for multiple comparisons.

\subsection{Approximate Firm-Size Elasticities}

GEIH reports size categories rather than exact worker
counts. To obtain an approximate elasticity, we assign sizes of 1, 2.5,
4.5, 8, 15, 25, 40.5, 75.5, and 150 to the nine firm-size categories in
their stated order. Let \(\widetilde{Size}_{iosdt}\) denote this assigned size.
We replace the size indicators with its logarithm and an interaction with
formality. Write \(\eta\) for the informal elasticity and \(\kappa\) for
the formal-minus-informal elasticity difference. The model is
\begin{equation}\label{eq:preferred-elasticity}
  \ln(w_{iosdt})=\alpha+\phi F_i
  +(\eta+\kappa F_i)\ln(\widetilde{Size}_{iosdt})
  +X_i'\gamma+\rho_o+\lambda_{s\times d\times t}
  +\varepsilon_{iosdt}.
\end{equation}
The formal elasticity is \(\eta+\kappa\). Both elasticities come from a single interacted regression,
with the preferred sample, controls, weights, and clustering. We test
\(H_0:\kappa=0\) to assess whether the slopes differ between formality groups. We also
report a raw regression of log income on log assigned size alone and a
common-slope version imposing
\(\kappa=0\), with formality and the preferred controls included in the latter.

These slopes depend on the assigned sizes, especially the value
assigned to 101+. They also impose a constant slope within each formality
group, whereas Equation~(\ref{eq:main}) allows each category to differ.
We compare them with the literature means reported by
\citet{DiegmannMullerSchoefer2026}, recognizing differences in populations and
controls. We assess the elasticity tests using the same alternative
clustering and bootstrap procedures. The appendix also
retains a version with fewer controls in Equation~(\ref{eq:elasticity}) for
comparison. We also estimate a quadratic extension, adding the square of log
assigned size and allowing both size coefficients to differ by formality.
Equation~(\ref{eq:polynomial-size-profile}) sets out that model. Its tests ask
whether a linear and a squared term per formality group can reproduce the
eight category-specific contrasts in Equation~(\ref{eq:main}), and check
sensitivity to the size assigned to 101+.

\subsection{Income Counterfactuals}
\label{subsec:income-counterfactual-method}

We construct illustrative income counterfactuals for the 341,282 observations
in the 2025 sample using the pooled preferred coefficients from
Equation~(\ref{eq:main}). Each scenario changes only firm size, formality,
or both. Worker characteristics, occupational position, sector, department,
weekly hours, and survey weights remain fixed. Workers already in the
destination category retain their baseline predictions.

Let \(\widehat{\ell}_i\) be fitted log hourly income from
Equation~(\ref{eq:main}), including the fixed effects, for worker \(i\)'s
observed size and formality status. Define \(S\) as the survey-weighted mean of
exponentiated residuals in 2025. We construct baseline hourly income as
\(\widehat w_i=S\exp(\widehat{\ell}_i)\). We assume the same retransformation
factor across workers and destinations. Let \(g(i)\) denote worker \(i\)'s
size category. If scenario
\(c\) changes the worker's estimated size-formality component
\(\widehat{\phi}F_i+\widehat{\beta}_{g(i)}+\widehat{\theta}_{g(i)}F_i\) by
\(\widehat\Delta_i^c\), counterfactual hourly income is
\(\widehat w_i^c=\widehat w_i\exp(\widehat\Delta_i^c)\).
For a change in size or formality alone, this estimated difference is obtained
from the corresponding contrast in Equation~(\ref{eq:preferred-contrasts}). For a joint change, we
subtract the origin component from the destination component directly. Let
\(\mathcal I_{2025}\) denote the 2025 sample. With expansion weight
\(a_i\) and weekly hours \(h_i\), the percentage change in total
predicted weekly labor income under scenario \(c\) is
\begin{equation}\label{eq:income-counterfactual}
  G_c=100\left[
  \frac{\sum_{i\in\mathcal I_{2025}}a_i h_i\widehat w_i
        \exp(\widehat\Delta_i^c)}
       {\sum_{i\in\mathcal I_{2025}}a_i h_i\widehat w_i}-1\right].
\end{equation}
This percentage change weights each individual's percentage change by that
individual's share of baseline income, including survey weights and hours.
It is not the average percentage wage change across workers. The common
factor \(S\) cancels from this ratio. The target-group calculation restricts
both sums to workers initially in the relevant formality group, including
those already at the destination. Restricting both sums to workers whose
size or formality changes gives the result for reassigned workers. Because
the workers and their weights remain the same before and after reassignment,
this is also the percentage change in their mean predicted weekly income.

\section{Results}

\subsection{From Raw to Conditional Wage Gaps}
\label{subsec:sequential-results}

Table~\ref{tab:firm-size-wage-premium-regressions} begins with the simplest
specification and adds controls in successive columns. Column (4) corresponds
to Equation~(\ref{eq:common-premium}). In column (1), which
contains only firm-size indicators, the 101+/solo coefficient is 1.017 log
points. Adding sector-year fixed effects lowers it to 0.817 in column (2).
Adding sex, age, age squared, and education lowers it to 0.598 in column (3),
and adding formality lowers it to 0.359 in column (4). The final coefficient
corresponds to a conditional premium of 43.2 percent. All four estimates are
statistically significant, and all four columns use the same sample.
The reduction therefore reflects the added controls rather than changes in
which observations enter the regressions.

\begin{table}[htbp]
  \centering
  \caption{Common-premium comparison: sequential controls}
  \label{tab:firm-size-wage-premium-regressions}
  \small
  \begin{tabular}{lcccc}
    \toprule
    & \multicolumn{4}{c}{Dependent variable: log real hourly labor income} \\
    \cmidrule(lr){2-5}
    & (1) & (2) & (3) & (4) \\
    \midrule
    Firm size: 2-3 & 0.184*** & 0.242*** & 0.209*** & 0.188*** \\
     & (0.034) & (0.031) & (0.027) & (0.029) \\
    Firm size: 4-5 & 0.363*** & 0.408*** & 0.361*** & 0.307*** \\
     & (0.034) & (0.030) & (0.029) & (0.038) \\
    Firm size: 6-10 & 0.499*** & 0.500*** & 0.422*** & 0.316*** \\
     & (0.029) & (0.033) & (0.037) & (0.050) \\
    Firm size: 11-19 & 0.625*** & 0.575*** & 0.457*** & 0.302*** \\
     & (0.034) & (0.040) & (0.042) & (0.050) \\
    Firm size: 20-30 & 0.687*** & 0.603*** & 0.462*** & 0.271*** \\
     & (0.046) & (0.045) & (0.046) & (0.051) \\
    Firm size: 31-50 & 0.759*** & 0.650*** & 0.487*** & 0.271*** \\
     & (0.056) & (0.049) & (0.047) & (0.048) \\
    Firm size: 51-100 & 0.818*** & 0.698*** & 0.518*** & 0.289*** \\
     & (0.063) & (0.051) & (0.048) & (0.047) \\
    Firm size: 101+ & 1.017*** & 0.817*** & 0.598*** & 0.359*** \\
     & (0.102) & (0.067) & (0.061) & (0.059) \\
    \midrule
    Gender, age, and education controls & No & No & Yes & Yes \\
    Formality control & No & No & No & Yes \\
    Sector-year fixed effects & No & Yes & Yes & Yes \\
    Observations & 5,268,514 & 5,268,514 & 5,268,514 & 5,268,514 \\
    $R^2$ & 0.213 & 0.271 & 0.376 & 0.386 \\
    \bottomrule
  \end{tabular}
  \vspace{0.3em}
  \begin{minipage}{0.95\textwidth}
  \footnotesize
  Notes: The omitted category is solo workers. All columns use the same estimation sample and GEIH expansion weights. Standard errors, clustered by sector, are reported in parentheses. Worker controls include a female-worker indicator, age, age squared, and education dummies. Column (4) adds labor formality and corresponds to Equation~(\ref{eq:common-premium}). All columns constrain size coefficients to be common across formality groups. Significance levels: * $p<0.10$, ** $p<0.05$, *** $p<0.01$.
  \end{minipage}
\end{table}


Observed worker characteristics, sector-year conditions, and formality account
for a substantial part of the raw association between size and pay. A positive
conditional gap remains. This pattern is consistent with evidence that worker
composition accounts for part, but not all, of size-related wage differences
\citep{BrownMedoff1989,troske1999evidence,Schaffner1998,
SoderbomTealWambugu2005,reed2019large,iregui2012diferenciales}.
The model still constrains the size coefficients to be common across formality
groups. Adding formality as a control allows a level difference between formal
and informal workers, but does not allow their firm-size premiums to differ.

Table~\ref{tab:firm-size-local-geography-robustness} keeps that restriction
and adds geographic and occupational-position controls. The 101+/solo
coefficient is 0.359 with sector-year effects, 0.338 after adding department-year
effects, and 0.325 with sector-department-year effects. The specification in
Equation~(\ref{eq:common-premium-geography}) also controls for occupational
position, giving a coefficient of 0.338 in
column (4), equivalent to a 40.2 percent premium. Thus, more detailed
controls do not eliminate the common-premium estimate, although it does not
decline at every step. The next subsection asks whether a common size premium
adequately describes formal and informal employment.

\begin{table}[htbp]
  \centering
  \caption{Common-premium estimates with geographic and occupational controls}
  \label{tab:firm-size-local-geography-robustness}
  \small
  \begin{tabular}{lcccc}
    \toprule
    & \multicolumn{4}{c}{Dependent variable: log real hourly labor income} \\
    \cmidrule(lr){2-5}
    & (1) & (2) & (3) & (4) \\
    \midrule
    Firm size: 2-3 & 0.188*** & 0.168*** & 0.167*** & 0.102*** \\
     & (0.011) & (0.011) & (0.011) & (0.013) \\
    Firm size: 4-5 & 0.307*** & 0.277*** & 0.273*** & 0.225*** \\
     & (0.014) & (0.014) & (0.014) & (0.015) \\
    Firm size: 6-10 & 0.316*** & 0.282*** & 0.278*** & 0.265*** \\
     & (0.018) & (0.017) & (0.017) & (0.016) \\
    Firm size: 11-19 & 0.302*** & 0.272*** & 0.267*** & 0.272*** \\
     & (0.018) & (0.017) & (0.017) & (0.015) \\
    Firm size: 20-30 & 0.271*** & 0.241*** & 0.237*** & 0.258*** \\
     & (0.019) & (0.017) & (0.017) & (0.015) \\
    Firm size: 31-50 & 0.271*** & 0.245*** & 0.240*** & 0.263*** \\
     & (0.018) & (0.017) & (0.017) & (0.015) \\
    Firm size: 51-100 & 0.289*** & 0.262*** & 0.257*** & 0.282*** \\
     & (0.018) & (0.017) & (0.017) & (0.015) \\
    Firm size: 101+ & 0.359*** & 0.338*** & 0.325*** & 0.338*** \\
     & (0.023) & (0.024) & (0.024) & (0.022) \\
    \midrule
    Worker and formality controls & Yes & Yes & Yes & Yes \\
    Occupational-position controls & No & No & No & Yes \\
    Fixed effects & $s\times t$ & $s\times t$, $d\times t$ & $s\times d\times t$ & $s\times d\times t$ \\
    Observations & 5,268,514 & 5,268,514 & 5,268,375 & 5,268,375 \\
    $R^2$ & 0.386 & 0.404 & 0.414 & 0.429 \\
    \bottomrule
  \end{tabular}
  \vspace{0.3em}
  \begin{minipage}{0.95\textwidth}
  \footnotesize
  Notes: The omitted category is solo workers. Column (4) corresponds to Equation~(\ref{eq:common-premium-geography}). $s$, $d$, and $t$ denote sector, department, and year. Parentheses contain sector-department-clustered standard errors. * $p<0.10$, ** $p<0.05$, *** $p<0.01$.
  \end{minipage}
\end{table}

\FloatBarrier

\subsection{Preferred Firm-Size Premiums by Formality}
\label{subsec:preferred-results}

Table~\ref{tab:firm-size-local-geography-formality-robustness} relaxes the
common-premium restriction by interacting each size indicator with formality.
Its four columns use the same control sequence as
Table~\ref{tab:firm-size-local-geography-robustness}. Column (4), our preferred model in
Equation~(\ref{eq:main}), includes worker characteristics, occupational-position
controls, and sector-department-year fixed effects. Columns (1)--(3) show how
the estimates change as geographic controls are added before controlling for
occupational position. All four columns allow the size coefficients to differ
between formal and informal workers and cluster standard errors by
sector-department cells.

\begin{table}[htbp]
  \centering
  \caption{Firm-size wage premiums by formality}
  \label{tab:firm-size-local-geography-formality-robustness}
  \scriptsize
  \begin{tabular}{lcccc}
    \toprule
    & \multicolumn{4}{c}{Dependent variable: log real hourly labor income} \\
    \cmidrule(lr){2-5}
    & (1) & (2) & (3) & \textbf{(4)} \\
    \midrule
    \multicolumn{5}{l}{\textit{Informal workers}} \\
    Firm size: 2-3 & 0.197*** {\scriptsize (0.010)} & 0.176*** {\scriptsize (0.010)} & 0.176*** {\scriptsize (0.010)} & 0.113*** {\scriptsize (0.013)} \\
    Firm size: 4-5 & 0.343*** {\scriptsize (0.012)} & 0.311*** {\scriptsize (0.011)} & 0.308*** {\scriptsize (0.011)} & 0.255*** {\scriptsize (0.016)} \\
    Firm size: 6-10 & 0.392*** {\scriptsize (0.013)} & 0.352*** {\scriptsize (0.012)} & 0.349*** {\scriptsize (0.012)} & 0.324*** {\scriptsize (0.017)} \\
    Firm size: 11-19 & 0.408*** {\scriptsize (0.016)} & 0.368*** {\scriptsize (0.014)} & 0.364*** {\scriptsize (0.014)} & 0.354*** {\scriptsize (0.018)} \\
    Firm size: 20-30 & 0.401*** {\scriptsize (0.016)} & 0.359*** {\scriptsize (0.013)} & 0.354*** {\scriptsize (0.014)} & 0.362*** {\scriptsize (0.017)} \\
    Firm size: 31-50 & 0.399*** {\scriptsize (0.018)} & 0.359*** {\scriptsize (0.018)} & 0.351*** {\scriptsize (0.018)} & 0.363*** {\scriptsize (0.020)} \\
    Firm size: 51-100 & 0.409*** {\scriptsize (0.020)} & 0.369*** {\scriptsize (0.019)} & 0.359*** {\scriptsize (0.019)} & 0.372*** {\scriptsize (0.020)} \\
    Firm size: 101+ & 0.336*** {\scriptsize (0.020)} & 0.304*** {\scriptsize (0.020)} & 0.288*** {\scriptsize (0.019)} & 0.316*** {\scriptsize (0.019)} \\
    \midrule
    \multicolumn{5}{l}{\textit{Formal workers}} \\
    Firm size: 2-3 & 0.022 {\scriptsize (0.029)} & 0.024 {\scriptsize (0.029)} & 0.021 {\scriptsize (0.029)} & -0.072** {\scriptsize (0.029)} \\
    Firm size: 4-5 & 0.025 {\scriptsize (0.037)} & 0.027 {\scriptsize (0.037)} & 0.023 {\scriptsize (0.037)} & -0.022 {\scriptsize (0.033)} \\
    Firm size: 6-10 & 0.023 {\scriptsize (0.036)} & 0.023 {\scriptsize (0.037)} & 0.019 {\scriptsize (0.037)} & 0.013 {\scriptsize (0.034)} \\
    Firm size: 11-19 & 0.035 {\scriptsize (0.034)} & 0.038 {\scriptsize (0.035)} & 0.033 {\scriptsize (0.034)} & 0.039 {\scriptsize (0.032)} \\
    Firm size: 20-30 & 0.024 {\scriptsize (0.034)} & 0.025 {\scriptsize (0.034)} & 0.021 {\scriptsize (0.034)} & 0.039 {\scriptsize (0.031)} \\
    Firm size: 31-50 & 0.042 {\scriptsize (0.033)} & 0.045 {\scriptsize (0.033)} & 0.040 {\scriptsize (0.033)} & 0.059* {\scriptsize (0.031)} \\
    Firm size: 51-100 & 0.069** {\scriptsize (0.032)} & 0.071** {\scriptsize (0.033)} & 0.066** {\scriptsize (0.033)} & 0.085*** {\scriptsize (0.031)} \\
    Firm size: 101+ & 0.154*** {\scriptsize (0.038)} & 0.161*** {\scriptsize (0.040)} & 0.147*** {\scriptsize (0.039)} & 0.153*** {\scriptsize (0.037)} \\
    \midrule
    Worker controls & Yes & Yes & Yes & Yes \\
    Occupational-position controls & No & No & No & Yes \\
    Fixed effects & $s\times t$ & $s\times t$, $d\times t$ & $s\times d\times t$ & $s\times d\times t$ \\
    Observations & 5,268,514 & 5,268,514 & 5,268,375 & 5,268,375 \\
    $R^2$ & 0.388 & 0.406 & 0.416 & 0.431 \\
    \bottomrule
  \end{tabular}
  \vspace{0.3em}
  \begin{minipage}{0.95\textwidth}
  \footnotesize
  Notes: Within-formality log-wage contrasts relative to solo workers. Each column comes from one interacted regression. Column (4) corresponds to Equation~(\ref{eq:main}). $s$, $d$, and $t$ denote sector, department, and year. Parentheses contain sector-department-clustered standard errors. * $p<0.10$, ** $p<0.05$, *** $p<0.01$.
  \end{minipage}
\end{table}


Among informal workers, column (4) shows a large wage gap between solo work and
employment in larger units. The premium relative to informal solo workers is
12.0 percent in firms with 2--3 workers, 29.0 percent in firms with 4--5 workers,
and 38.2 percent in firms with 6--10 workers. It reaches 42.4--45.1 percent
across the 11--100 categories and is 37.1 percent in 101+ firms. All these
comparisons with informal solo work are statistically significant. The premium
therefore rises sharply across the smallest categories, but does not increase
monotonically through the largest category.

Among formal workers, the preferred estimates also reveal a large-firm
advantage in the point estimates. Formal workers in 101+ firms earn an estimated
16.5 percent more than observably comparable formal solo workers. The 51--100
category has an 8.9 percent premium over formal solo work. The smaller categories
do not form a uniformly increasing sequence relative to solo work. In
particular, formal workers in the 2--3 category earn 7.0 percent less than formal solo workers.
The estimates for 4--50 workers range from -2.2 to 6.1 percent.

Figure~\ref{fig:preferred-formality-premium} displays the within-formality
contrasts from Equation~(\ref{eq:main}). The larger informal premium relative to solo work
does not imply that every increase in firm size corresponds to higher wages.
The informal estimates flatten across intermediate categories and fall at
101+, whereas the formal 101+ estimate is the largest in its panel.

\begin{figure}[htbp]
  \centering
  \includegraphics[width=0.95\textwidth]{preferred_formality_premiums_\figurelang.png}
  \caption{Firm-size premiums by formality in the preferred specification}
  \label{fig:preferred-formality-premium}
  \begin{minipage}{0.95\textwidth}
    \footnotesize
    Notes: Confidence intervals use sector-department clustering.
  \end{minipage}
\end{figure}

Appendix Table~\ref{tab:preferred-clustering-sensitivity} compares the
specifications in columns (1), (3), and (4) of
Table~\ref{tab:firm-size-local-geography-formality-robustness} on a common
sample under both clustering conventions. The formal 101+/solo estimate
is 16.6 percent with sector-year effects and 16.5 percent with the preferred
controls. The underlying estimated log-wage contrast,
\(\widehat{\Delta}^{F}_{g_0\to101+}\), has a standard error of 0.037 log points
under sector-department clustering, compared with 0.091 under sector
clustering. The corresponding two-sided \(p\)-values are below 0.001 and
0.113. Appendix Tables~\ref{tab:preferred-clustering-comparison} and
\ref{tab:preferred-wild-bootstrap} extend this check to department and
two-way clustering and to wild cluster bootstrap inference. The formal
101+/solo bootstrap \(p\)-value is 0.184 with sector clustering and 0.188
with two-way clustering. None of the formal size-to-solo contrasts is
significant at 5 percent under either bootstrap convention. All eight
informal size-to-solo contrasts remain significant. The preferred regression
tables in the main text and their associated confidence intervals retain
sector-department clustering, but the formal
101+/solo significance should not be treated as robust.

Table~\ref{tab:firm-size-local-geography-formality-robustness-101-reference}
reports the specifications in
Table~\ref{tab:firm-size-local-geography-formality-robustness} with 101+ workers
as the reference within each formality group. In column (4), which corresponds
to Equation~(\ref{eq:main}), every smaller formal category has a lower
point estimate than formal 101+ workers. All seven comparisons with formal
categories from 2--3 through 51--100 remain significant at 5 percent under
the alternative clustering and bootstrap procedures, unlike the solo comparison.
Expressed with the smaller category as the denominator, the 101+ premium is
approximately 7.0 percent relative to 51--100 workers and 25.2 percent relative
to 2--3 workers. Among informal workers, by contrast, the point estimates
put 101+ workers above solo workers and workers in firms with 2--5 workers, but not above
every intermediate category.

\begin{table}[htbp]
  \centering
  \caption{Firm-size wage premiums by formality: reference category 101+}
  \label{tab:firm-size-local-geography-formality-robustness-101-reference}
  \scriptsize
  \begin{tabular}{lcccc}
    \toprule
    & \multicolumn{4}{c}{Dependent variable: log real hourly labor income} \\
    \cmidrule(lr){2-5}
    & (1) & (2) & (3) & \textbf{(4)} \\
    \midrule
    \multicolumn{5}{l}{\textit{Informal workers}} \\
    Firm size: Solo & -0.341*** {\scriptsize (0.020)} & -0.309*** {\scriptsize (0.020)} & -0.292*** {\scriptsize (0.020)} & -0.315*** {\scriptsize (0.019)} \\
    Firm size: 2-3 & -0.138*** {\scriptsize (0.018)} & -0.127*** {\scriptsize (0.017)} & -0.110*** {\scriptsize (0.016)} & -0.203*** {\scriptsize (0.016)} \\
    Firm size: 4-5 & 0.011 {\scriptsize (0.016)} & 0.011 {\scriptsize (0.016)} & 0.024 {\scriptsize (0.016)} & -0.058*** {\scriptsize (0.017)} \\
    Firm size: 6-10 & 0.058*** {\scriptsize (0.016)} & 0.050*** {\scriptsize (0.016)} & 0.062*** {\scriptsize (0.016)} & 0.009 {\scriptsize (0.017)} \\
    Firm size: 11-19 & 0.074*** {\scriptsize (0.016)} & 0.066*** {\scriptsize (0.016)} & 0.077*** {\scriptsize (0.016)} & 0.040** {\scriptsize (0.017)} \\
    Firm size: 20-30 & 0.069*** {\scriptsize (0.016)} & 0.059*** {\scriptsize (0.016)} & 0.070*** {\scriptsize (0.016)} & 0.050*** {\scriptsize (0.016)} \\
    Firm size: 31-50 & 0.065*** {\scriptsize (0.016)} & 0.057*** {\scriptsize (0.016)} & 0.065*** {\scriptsize (0.016)} & 0.048*** {\scriptsize (0.016)} \\
    Firm size: 51-100 & 0.078*** {\scriptsize (0.016)} & 0.068*** {\scriptsize (0.016)} & 0.074*** {\scriptsize (0.016)} & 0.060*** {\scriptsize (0.016)} \\
    \midrule
    \multicolumn{5}{l}{\textit{Formal workers}} \\
    Firm size: Solo & -0.167*** {\scriptsize (0.038)} & -0.174*** {\scriptsize (0.039)} & -0.160*** {\scriptsize (0.039)} & -0.159*** {\scriptsize (0.037)} \\
    Firm size: 2-3 & -0.145*** {\scriptsize (0.018)} & -0.149*** {\scriptsize (0.018)} & -0.138*** {\scriptsize (0.018)} & -0.234*** {\scriptsize (0.018)} \\
    Firm size: 4-5 & -0.137*** {\scriptsize (0.012)} & -0.142*** {\scriptsize (0.013)} & -0.132*** {\scriptsize (0.012)} & -0.182*** {\scriptsize (0.012)} \\
    Firm size: 6-10 & -0.138*** {\scriptsize (0.012)} & -0.144*** {\scriptsize (0.013)} & -0.134*** {\scriptsize (0.012)} & -0.146*** {\scriptsize (0.011)} \\
    Firm size: 11-19 & -0.124*** {\scriptsize (0.015)} & -0.129*** {\scriptsize (0.016)} & -0.119*** {\scriptsize (0.016)} & -0.118*** {\scriptsize (0.014)} \\
    Firm size: 20-30 & -0.133*** {\scriptsize (0.014)} & -0.139*** {\scriptsize (0.015)} & -0.129*** {\scriptsize (0.015)} & -0.116*** {\scriptsize (0.012)} \\
    Firm size: 31-50 & -0.114*** {\scriptsize (0.013)} & -0.118*** {\scriptsize (0.014)} & -0.109*** {\scriptsize (0.014)} & -0.095*** {\scriptsize (0.012)} \\
    Firm size: 51-100 & -0.088*** {\scriptsize (0.012)} & -0.093*** {\scriptsize (0.013)} & -0.084*** {\scriptsize (0.013)} & -0.070*** {\scriptsize (0.011)} \\
    \midrule
    Worker controls & Yes & Yes & Yes & Yes \\
    Occupational-position controls & No & No & No & Yes \\
    Fixed effects & $s\times t$ & $s\times t$, $d\times t$ & $s\times d\times t$ & $s\times d\times t$ \\
    Observations & 5,426,316 & 5,426,316 & 5,426,174 & 5,426,174 \\
    $R^2$ & 0.394 & 0.412 & 0.422 & 0.437 \\
    \bottomrule
  \end{tabular}
  \vspace{0.3em}
  \begin{minipage}{0.95\textwidth}
  \footnotesize
  Notes: Within-formality log-wage contrasts relative to 101+ workers, from the specifications in Table~\ref{tab:firm-size-local-geography-formality-robustness}. $s$, $d$, and $t$ denote sector, department, and year. Parentheses contain sector-department-clustered standard errors. * $p<0.10$, ** $p<0.05$, *** $p<0.01$.
  \end{minipage}
\end{table}


Table~\ref{tab:firm-size-local-geography-formality-robustness-11-19-reference}
makes the intermediate comparison explicit. In column (4), formal 101+ workers
earn 0.113 log points, or approximately 12.0 percent, more than formal workers
in firms with 11--19 workers. Informal 101+ workers, however, earn 0.038 log
points, or approximately 3.7 percent, less than informal workers in that
reference category. The formal contrast remains significant under two-way
bootstrap inference (\(p=0.003\)), whereas the informal contrast does not
(\(p=0.210\)), despite its significance under sector-department clustering.
The same regression therefore implies different answers
depending on both formality and the sizes being compared. These are
alternative normalizations of each specification, not separate robustness
checks. In particular, column (4) in all three tables comes from the same
preferred regression in Equation~(\ref{eq:main}).

\begin{table}[p]
  \centering
  \caption{Firm-size wage premiums by formality: reference category 11--19}
  \label{tab:firm-size-local-geography-formality-robustness-11-19-reference}
  \scriptsize
  \begin{tabular}{lcccc}
    \toprule
    & \multicolumn{4}{c}{Dependent variable: log real hourly labor income} \\
    \cmidrule(lr){2-5}
    & (1) & (2) & (3) & \textbf{(4)} \\
    \midrule
    \multicolumn{5}{l}{\textit{Informal workers}} \\
    Firm size: Solo & -0.408*** {\scriptsize (0.016)} & -0.368*** {\scriptsize (0.014)} & -0.364*** {\scriptsize (0.014)} & -0.354*** {\scriptsize (0.018)} \\
    Firm size: 2-3 & -0.211*** {\scriptsize (0.013)} & -0.192*** {\scriptsize (0.011)} & -0.187*** {\scriptsize (0.011)} & -0.240*** {\scriptsize (0.011)} \\
    Firm size: 4-5 & -0.065*** {\scriptsize (0.009)} & -0.058*** {\scriptsize (0.009)} & -0.056*** {\scriptsize (0.009)} & -0.099*** {\scriptsize (0.007)} \\
    Firm size: 6-10 & -0.016*** {\scriptsize (0.006)} & -0.016*** {\scriptsize (0.006)} & -0.015** {\scriptsize (0.006)} & -0.030*** {\scriptsize (0.006)} \\
    Firm size: 20-30 & -0.007 {\scriptsize (0.008)} & -0.010 {\scriptsize (0.007)} & -0.009 {\scriptsize (0.007)} & 0.008 {\scriptsize (0.007)} \\
    Firm size: 31-50 & -0.009 {\scriptsize (0.011)} & -0.009 {\scriptsize (0.011)} & -0.012 {\scriptsize (0.011)} & 0.009 {\scriptsize (0.010)} \\
    Firm size: 51-100 & 0.001 {\scriptsize (0.014)} & 0.000 {\scriptsize (0.013)} & -0.004 {\scriptsize (0.014)} & 0.019 {\scriptsize (0.013)} \\
    Firm size: 101+ & -0.072*** {\scriptsize (0.016)} & -0.065*** {\scriptsize (0.016)} & -0.075*** {\scriptsize (0.016)} & -0.038** {\scriptsize (0.017)} \\
    \midrule
    \multicolumn{5}{l}{\textit{Formal workers}} \\
    Firm size: Solo & -0.035 {\scriptsize (0.034)} & -0.038 {\scriptsize (0.035)} & -0.033 {\scriptsize (0.034)} & -0.039 {\scriptsize (0.032)} \\
    Firm size: 2-3 & -0.013 {\scriptsize (0.014)} & -0.014 {\scriptsize (0.014)} & -0.013 {\scriptsize (0.013)} & -0.112*** {\scriptsize (0.014)} \\
    Firm size: 4-5 & -0.010 {\scriptsize (0.009)} & -0.011 {\scriptsize (0.008)} & -0.010 {\scriptsize (0.008)} & -0.062*** {\scriptsize (0.008)} \\
    Firm size: 6-10 & -0.012* {\scriptsize (0.007)} & -0.015** {\scriptsize (0.007)} & -0.014** {\scriptsize (0.007)} & -0.027*** {\scriptsize (0.006)} \\
    Firm size: 20-30 & -0.011** {\scriptsize (0.005)} & -0.013*** {\scriptsize (0.004)} & -0.012*** {\scriptsize (0.004)} & -0.000 {\scriptsize (0.004)} \\
    Firm size: 31-50 & 0.007 {\scriptsize (0.005)} & 0.008 {\scriptsize (0.005)} & 0.007 {\scriptsize (0.005)} & 0.020*** {\scriptsize (0.005)} \\
    Firm size: 51-100 & 0.034*** {\scriptsize (0.007)} & 0.033*** {\scriptsize (0.007)} & 0.033*** {\scriptsize (0.007)} & 0.046*** {\scriptsize (0.007)} \\
    Firm size: 101+ & 0.119*** {\scriptsize (0.015)} & 0.123*** {\scriptsize (0.016)} & 0.114*** {\scriptsize (0.016)} & 0.113*** {\scriptsize (0.014)} \\
    \midrule
    Worker controls & Yes & Yes & Yes & Yes \\
    Occupational-position controls & No & No & No & Yes \\
    Fixed effects & $s\times t$ & $s\times t$, $d\times t$ & $s\times d\times t$ & $s\times d\times t$ \\
    Observations & 5,268,514 & 5,268,514 & 5,268,375 & 5,268,375 \\
    $R^2$ & 0.388 & 0.406 & 0.416 & 0.431 \\
    \bottomrule
  \end{tabular}
  \vspace{0.3em}
  \begin{minipage}{0.95\textwidth}
  \footnotesize
  Notes: The omitted categories are workers in firms with 11--19 workers within each formality group. Standard errors, clustered by sector-department cells, are reported in parentheses. Worker controls include a female-worker indicator, age, age squared, and education dummies. The local geography $d$ is the department, treating Bogot\'a as a department-equivalent Capital District. Column (4) is the preferred specification in Equation~(\ref{eq:main}). Significance levels: * $p<0.10$, ** $p<0.05$, *** $p<0.01$.
  \end{minipage}
\end{table}


Table~\ref{tab:preferred-formality-gap} reports formal-informal wage gaps
from Equation~(\ref{eq:main}), using the contrasts defined in
Equation~(\ref{eq:preferred-contrasts}). Formal workers earn 59.3 percent more
than informal workers within solo work, 32.4 percent more in 2--3 firms,
and 35.4 percent more in 101+ firms. Across the 4--100 categories, the gaps
range from 15.4 to 20.8 percent. All nine gaps remain significant under
each alternative clustering and bootstrap procedure. These comparisons hold firm size fixed, whereas
Figure~\ref{fig:preferred-formality-premium} holds formality fixed.

\begin{table}[htbp]
  \centering
  \caption{Formal-informal wage gaps in the preferred specification}
  \label{tab:preferred-formality-gap}
  \small
  \begin{tabular}{lccc}
    \toprule
    Firm size & Log-wage gap & Premium & $p$-value \\
    \midrule
    Solo & 0.468*** & 59.6\% & $<0.001$ \\
    2-3 & 0.281*** & 32.4\% & $<0.001$ \\
    4-5 & 0.188*** & 20.7\% & $<0.001$ \\
    6-10 & 0.157*** & 17.0\% & $<0.001$ \\
    11-19 & 0.155*** & 16.7\% & $<0.001$ \\
    20-30 & 0.146*** & 15.8\% & $<0.001$ \\
    31-50 & 0.169*** & 18.4\% & $<0.001$ \\
    51-100 & 0.182*** & 19.9\% & $<0.001$ \\
    101+ & 0.312*** & 36.6\% & $<0.001$ \\
    \bottomrule
  \end{tabular}
  \vspace{0.3em}
  \begin{minipage}{0.95\textwidth}
  \footnotesize
  Notes: Formal-minus-informal contrasts from Equation~(\ref{eq:main}) at each firm size. Percentage premiums use informal workers as the denominator. Two-sided tests use sector-department-clustered standard errors. * $p<0.10$, ** $p<0.05$, *** $p<0.01$.
  \end{minipage}
\end{table}


The contrast with the literature is most direct for the informal profile.
\citet{BalkanTumen2016} find a more pronounced firm-size wage gap among informal
than formal jobs in Turkey. Our preferred estimates show a similar contrast
when larger categories are compared with solo work. They also show a
statistically significant advantage of the largest formal employers over
the non-solo formal categories. This distinction complements evidence that firm
size and formality are closely related in Colombia
\citep{OteroCortesEtAl2025,DazaGamboa2013,GutierrezRodriguezLesmes2023} and the
Ecuadorian evidence that controlling for firm size can eliminate an apparent
formal-sector premium in wages net of taxes \citep{ElBadaouiStroblWalsh2010}. The Colombian
within-formality size comparisons cannot be reduced to a contrast between
formal and informal employment.

Appendix Figure~\ref{fig:firm-size-wage-premium} displays the common-premium
estimates from Equation~(\ref{eq:common-premium}), reported in column (4) of
Table~\ref{tab:firm-size-wage-premium-regressions}.
Appendix Tables~\ref{tab:firm-size-wage-premium-regressions-101-reference}
and \ref{tab:firm-size-wage-premium-regressions-11-19-reference} re-express
all four specifications in that table with 101+ and 11--19 workers as their
respective reference categories.
Tables~\ref{tab:firm-size-local-geography-robustness-101-reference} and
\ref{tab:firm-size-local-geography-robustness-11-19-reference} do the same for
the geographic and occupational-position specifications in
Table~\ref{tab:firm-size-local-geography-robustness}.

Appendix Table~\ref{tab:formality-interaction-coefficients} and
Figure~\ref{fig:firm-size-formality-premium} report the simpler interacted
model in Equation~(\ref{eq:formality-interaction}), with sector-year controls.
Table~\ref{tab:formal-101-contrasts} reports its
formal 101+ comparisons, and Table~\ref{tab:formal-wage-gap-by-size} its
formal-informal gaps. Tables~\ref{tab:formality-interaction-coefficients-101-reference}
and \ref{tab:formality-interaction-coefficients-11-19-reference} show its
alternative normalizations. These simpler results do not replace the
preferred comparisons above.

\subsection{Salaried Workers}
\label{subsec:salaried-workers}

Table~\ref{tab:salaried-robustness} compares estimates of
Equation~(\ref{eq:main}) for the full sample and for private-sector employees,
government employees, domestic workers, and day laborers only. The restricted
sample contains 2,682,009 observations.
The main rows use 11--19 workers as the reference within each formality group.
Relative to that category, the formal 101+ premium is 12.9 percent, compared
with 12.0 percent in the full sample. The informal 101+ gap is -3.2 percent,
compared with -3.7 percent. Both contrasts are significant at
5 percent under sector-department clustering.

\begin{table}[!htbp]
  \centering
  \caption{Firm-size premiums in the full and salaried-only samples}
  \label{tab:salaried-robustness}
  \small
  \begin{tabular}{lrrrr}
    \toprule
    & \multicolumn{2}{c}{All workers} & \multicolumn{2}{c}{Salaried workers} \\
    \cmidrule(lr){2-3}\cmidrule(lr){4-5}
    Firm size & Informal & Formal & Informal & Formal \\
    \midrule
    Solo & -0.354*** & -0.041 & -0.226*** & -0.064*** \\
     & (0.018) & (0.031) & (0.016) & (0.023) \\
    2--3 & -0.242*** & -0.116*** & -0.180*** & -0.082*** \\
     & (0.011) & (0.014) & (0.011) & (0.013) \\
    4--5 & -0.098*** & -0.064*** & -0.080*** & -0.058*** \\
     & (0.007) & (0.007) & (0.007) & (0.007) \\
    6--10 & -0.030*** & -0.028*** & -0.028*** & -0.022*** \\
     & (0.006) & (0.007) & (0.005) & (0.006) \\
    11--19 & Reference & Reference & Reference & Reference \\
    20--30 & 0.010 & 0.002 & 0.011* & 0.011*** \\
     & (0.007) & (0.004) & (0.007) & (0.003) \\
    31--50 & 0.009 & 0.023*** & 0.009 & 0.035*** \\
     & (0.010) & (0.005) & (0.010) & (0.004) \\
    51--100 & 0.021 & 0.048*** & 0.030** & 0.062*** \\
     & (0.013) & (0.006) & (0.012) & (0.006) \\
    101+ & -0.040** & 0.118*** & -0.033** & 0.125*** \\
     & (0.017) & (0.014) & (0.017) & (0.012) \\
    \midrule
    11--19/solo premium (\%) & 42.5 & 4.2 & 25.3 & 6.6 \\
    101+/solo premium (\%) & 37.0 & 17.2 & 21.3 & 20.8 \\
    Equal 101+/solo premiums: $p$ & \multicolumn{2}{c}{$<0.001$} & \multicolumn{2}{c}{0.888} \\
    Equal size profiles: $p$ & \multicolumn{2}{c}{$<0.001$} & \multicolumn{2}{c}{$<0.001$} \\
    Observations & \multicolumn{2}{c}{5,426,174} & \multicolumn{2}{c}{2,765,713} \\
    \bottomrule
  \end{tabular}
  \vspace{0.3em}
  \begin{minipage}{0.95\textwidth}
  \footnotesize
  Notes: Main rows report within-formality log-wage contrasts relative to 11--19 workers. Summary premiums use solo work as the reference and are in percent. Each column pair estimates Equation~(\ref{eq:main}) on the indicated sample. The restricted sample excludes 66 additional singleton cells. Standard errors in parentheses use sector-department clustering, with $G=559$ clusters in the full sample and $G=555$ among employees. Equality tests compare formal and informal log contrasts relative to solo work, individually for 101+ (two-sided $t$ test) and jointly across all sizes (Wald $F$ test with eight numerator degrees of freedom). Both use $G-1$ denominator degrees of freedom. * $p<0.10$, ** $p<0.05$, *** $p<0.01$.
  \end{minipage}
\end{table}

\FloatBarrier

Figure~\ref{fig:salaried-formality-premium} displays the premiums from
Equation~(\ref{eq:main}) estimated for salaried workers only, using the same
vertical scale as Figure~\ref{fig:preferred-formality-premium} and comparing
each size with solo workers of the same formality status.
The conditional 101+/solo premium is 20.2 percent
among formal employees and 20.6 percent among informal employees, and a direct
test does not reject equality (\(p=0.897\)). This similarity does not extend
to the entire size profile. Relative to solo work, the premium in firms with
11--19 workers remains substantially larger among informal employees, at
24.7 percent, than among formal employees, at 6.5 percent. The table's summary
rows report both comparisons. A Wald test rejects the joint null that all eight
size contrasts relative to solo work are equal across formality groups (\(p<0.001\)).
Both conclusions also hold under sector
clustering. Thus, the larger informal premium in the full sample's 101+/solo
comparison does not generalize to salaried employment alone, although
formality-related differences across intermediate firm sizes remain.

\begin{figure}[htbp]
  \centering
  \includegraphics[width=0.95\textwidth]{salaried_formality_premiums_\figurelang.png}
  \caption{Firm-size premiums by formality among salaried workers}
  \label{fig:salaried-formality-premium}
  \begin{minipage}{0.95\textwidth}
    \footnotesize
    Notes: Confidence intervals use sector-department clustering.
  \end{minipage}
\end{figure}
\FloatBarrier

\subsection{Dynamic Results}
\label{subsec:dynamic-results}

The dynamic exercise asks whether the conditional wage gap between a given
firm-size category and solo work widened or narrowed within each formality
group. We use Equation~(\ref{eq:preferred-dynamics}), allowing all worker and
occupational-position coefficients to vary by year while retaining
sector-department fixed effects within each year.

Figure~\ref{fig:preferred-formality-endpoints} compares the adjusted profiles
from Equation~(\ref{eq:preferred-dynamics}) estimated separately for 2008
and 2025. Among formal workers, the 101+/solo premium changes little,
from 12.6 to 13.5 percent. The 2--3 category remains below formal solo work,
at -11.9 percent in 2008 and -10.9 percent in 2025. The intermediate formal
categories also show relatively small endpoint differences. Among informal
workers, the 101+/solo premium rises from 35.6 to 57.7 percent. The premiums
for 2--3, 4--5, and 6--10 workers rise from 6.1, 23.7, and 28.2 percent
to 18.7, 37.4, and 47.4 percent, respectively.

\begin{figure}[htbp]
  \centering
  \includegraphics[width=0.95\textwidth]{preferred_formality_endpoints_\figurelang.png}
  \caption{Firm-size premiums by formality with preferred controls, 2008 and 2025}
  \label{fig:preferred-formality-endpoints}
  \begin{minipage}{0.95\textwidth}
    \footnotesize
    Notes: Confidence intervals use sector-department clustering.
  \end{minipage}
\end{figure}

Table~\ref{tab:preferred-endpoint-changes} tests the 2008--2025 change for each
size-formality comparison using Equation~(\ref{eq:preferred-dynamics})
estimated jointly for the two years.
The formal 101+/solo premium increases by
0.9 percentage points, with \(p=0.861\). None of the formal endpoint changes
is significant at 5 percent. The informal 101+/solo premium increases by
22.2 percentage points, with \(p=0.011\) under sector-department clustering.
The informal increases for 2--3, 4--5, and 6--10 workers are also significant.
The remaining informal changes are not significant at 5 percent in the
pointwise tests. Appendix Table~\ref{tab:preferred-clustering-comparison}
shows that sector, department, and two-way clustering raise the \(p\)-value
for the informal 101+/solo change to 0.090, 0.124, and 0.199, respectively.
The corresponding wild bootstrap \(p\)-values in
Table~\ref{tab:preferred-wild-bootstrap} are 0.111, 0.266, and 0.273.
Thus, the point estimates suggest a sizeable widening, but the increase
is not statistically significant under any of these bootstrap procedures.

\begin{table}[htbp]
  \centering
  \caption{Changes in within-formality firm-size premiums, 2008--2025}
  \label{tab:preferred-endpoint-changes}
  \small
  \begin{tabular}{lrrrr}
    \toprule
    Firm size & 2008 (\%) & 2025 (\%) & Change (pp) & $p$-value \\
    \midrule
    \multicolumn{5}{l}{\textit{Formal workers}} \\
    2-3 & -7.7 & -11.2 & -3.5 & 0.265 \\
    4-5 & -2.6 & -3.1 & -0.5 & 0.883 \\
    6-10 & 4.1 & -1.8 & -5.8 & 0.117 \\
    11-19 & 5.5 & 2.4 & -3.1 & 0.475 \\
    20-30 & 5.5 & 0.9 & -4.7 & 0.266 \\
    31-50 & 6.2 & 4.0 & -2.2 & 0.564 \\
    51-100 & 5.7 & 6.3 & 0.5 & 0.899 \\
    101+ & 16.6 & 13.5 & -3.2 & 0.467 \\
    \midrule
    \multicolumn{5}{l}{\textit{Informal workers}} \\
    2-3 & 7.5 & 18.2 & 10.7 & $<0.001$ \\
    4-5 & 23.3 & 36.8 & 13.5 & 0.002 \\
    6-10 & 29.4 & 47.2 & 17.8 & $<0.001$ \\
    11-19 & 37.2 & 49.6 & 12.4 & 0.034 \\
    20-30 & 44.5 & 44.5 & 0.0 & 0.997 \\
    31-50 & 39.8 & 42.5 & 2.7 & 0.679 \\
    51-100 & 43.6 & 47.3 & 3.7 & 0.639 \\
    101+ & 33.2 & 57.1 & 23.8 & 0.002 \\
    \bottomrule
  \end{tabular}
  \vspace{0.3em}
  \begin{minipage}{0.95\textwidth}
  \footnotesize
  Notes: Estimates from Equation~(\ref{eq:preferred-dynamics}), with solo work as the reference within each formality-year group. Changes are in percentage points and calculated before rounding. Two-sided $p$-values test equality of the endpoint log-wage contrasts using joint estimation and sector-department clustering. Tests are pointwise, not adjusted for multiple comparisons.
  \end{minipage}
\end{table}


Figure~\ref{fig:preferred-large-firm-dynamics} reports the annual 101+/solo
premiums from Equation~(\ref{eq:preferred-dynamics}), showing why the endpoint
comparison should not be read as a steady trend. The formal premium reaches
25.7 percent in 2012, falls to 5.9 percent in 2021, and recovers to
13.5 percent in 2025. With sector-department clustering, its confidence
interval excludes zero in every year except 2021, and the informal premium
is positive and significant throughout. The latter also fluctuates,
falling to 28.7 percent in 2021 before
reaching 57.7 percent in 2025.

\begin{figure}[htbp]
  \centering
  \includegraphics[width=0.95\textwidth]{preferred_large_firm_dynamics_\figurelang.png}
  \caption{Large-firm premiums by formality with preferred controls in each year}
  \label{fig:preferred-large-firm-dynamics}
  \begin{minipage}{0.95\textwidth}
    \footnotesize
    Notes: Confidence intervals use sector-department clustering.
  \end{minipage}
\end{figure}

These results qualify comparisons with the literature on the evolution of
the premium. \citet{BloomGuvenenSmithSongVonWachter2018} document a decline
in the United States, and \citet{reed2019large} find that the premium tends
to decline with national income and over time across developing economies.
\citet{ColonnelliTagWebbWolter2018} find no common trajectory across Brazil,
Germany, Sweden, and the United Kingdom. Our preferred-control estimates
show little net change between the formal endpoints and a larger informal
gap in 2025, but the informal widening is not significant under the
bootstrap checks.
They do not describe a uniform decline across Colombian workers.

For comparison, Equation~(\ref{eq:dynamic-benchmark}) allows each size
coefficient to vary by year while retaining common
size premiums across formality groups and constant worker-control coefficients.
Appendix Figure~\ref{fig:firm-size-wage-premium-over-time} reports these
annual estimates. The companion model in Equation~(\ref{eq:trend-benchmark})
imposes a constant annual change in each size-to-solo log-wage gap.
Table~\ref{tab:firm-size-premium-trend-test}
reports positive but statistically insignificant annual changes, with
alternative reference categories in
Tables~\ref{tab:firm-size-premium-trend-test-101-reference} and
\ref{tab:firm-size-premium-trend-test-11-19-reference}. Appendix
Figures~\ref{fig:adjusted-formality-endpoints} and
\ref{fig:large-firm-formality-premium-over-time} report the simpler
formality-specific profiles from Equation~(\ref{eq:formality-interaction})
estimated separately by year. These are comparison specifications, not the
source of the preferred dynamic results above.

\subsection{Benchmarking the Firm-Size Wage Elasticity}
\label{subsec:elasticity-benchmark}

Table~\ref{tab:preferred-elasticities} compares the approximate elasticities
from Equation~(\ref{eq:preferred-elasticity}) and its raw and common-slope
variants with the no-controls and
observable-controls means reported by \citet{DiegmannMullerSchoefer2026}.
The raw Colombian elasticity is 0.203, compared with their no-controls
mean of 0.062. After worker characteristics, occupational position,
sector-department-year effects, and formality are added, a model imposing
a common slope gives 0.058, compared with the literature mean of 0.036.

\begin{table}[htbp]
  \centering
  \caption{Approximate firm-size elasticities with preferred controls}
  \label{tab:preferred-elasticities}
  \small
  \begin{tabular}{p{0.29\textwidth}p{0.37\textwidth}cc}
    \toprule
    Source & Specification & Elasticity & S.E. \\
    \midrule
    Diegmann et al. (2026) & No worker controls & 0.062 & n.a. \\
    Diegmann et al. (2026) & Observable worker controls & 0.036 & n.a. \\
    \midrule
    Colombia (GEIH) & No controls & 0.203*** & (0.008) \\
    Colombia (GEIH) & Common-slope model & 0.058*** & (0.004) \\
    Colombia (GEIH) & Formal workers & 0.040*** & (0.005) \\
    Colombia (GEIH) & Informal workers & 0.097*** & (0.004) \\
    \bottomrule
  \end{tabular}
  \vspace{0.3em}
  \begin{minipage}{0.95\textwidth}
  \footnotesize
  Notes: Literature means reproduce Table 1 of \citet{DiegmannMullerSchoefer2026}. All Colombian estimates use the preferred sample and GEIH weights. Controlled models include sex, age, age squared, education, occupational position, and sector-department-year fixed effects. The common-slope model also includes formality. Formal and informal slopes come from a single interaction model, Equation~(\ref{eq:preferred-elasticity}). Assigned sizes are 1, 2.5, 4.5, 8, 15, 25, 40.5, 75.5, and 150. Standard errors are clustered by sector-department. * $p<0.10$, ** $p<0.05$, *** $p<0.01$.
  \end{minipage}
\end{table}


Allowing the slope to differ by formality gives an elasticity of 0.040 for
formal workers and 0.097 for informal workers. The formal estimate is close
in magnitude to the literature's observable-controls mean, whereas the informal
estimate is about 2.7 times that benchmark. The informal slope exceeds the
formal slope by \(-\widehat{\kappa}=0.057\), a statistically significant difference
(two-way bootstrap \(p=0.0015\)). These are two contrasts from one regression, with common
worker and occupational-position coefficients. They are not estimates from
separate samples or a decomposition of the common elasticity.
Appendix Tables~\ref{tab:preferred-clustering-comparison} and
\ref{tab:preferred-wild-bootstrap} show that both slopes and their difference
remain significant under all the clustering and bootstrap procedures examined.

The literature comparison concerns magnitudes across different populations,
size measures, and controls, not tests against a common population parameter.
The assigned-size slopes also summarize the entire size distribution.
A positive informal slope does not contradict the lower categorical estimate
for informal 101+ workers than for several intermediate categories.
The categorical preferred model does not impose the elasticity model's
constant slope within each formality group.

Appendix Table~\ref{tab:firm-size-elasticity-benchmark} retains the simpler
elasticities of 0.026 for formal workers and 0.130 for informal workers.
They come from separate-formality regressions based on Equation~(\ref{eq:elasticity}),
with sector-year controls and no occupational-position controls.
Appendix Table~\ref{tab:polynomial-size-profile} reports the quadratic model in
Equation~(\ref{eq:polynomial-size-profile}), which uses the preferred controls.
Figure~\ref{fig:polynomial-size-profile} overlays its curves on the categorical
estimates from Equation~(\ref{eq:main}).
The curves summarize the broad profiles but do not reproduce all categorical
differences. The restriction tests in
Table~\ref{tab:quadratic-preferred-sensitivity} reject this restriction on
Equation~(\ref{eq:main}) within
both formality groups, including under alternative assignments for 101+.
The quadratic is therefore a descriptive approximation, not a replacement
for the preferred categorical estimates.

\FloatBarrier
\subsection{Illustrative Income Counterfactuals}
\label{subsec:income-counterfactuals}

Table~\ref{tab:income-counterfactuals-2025} reports four scenarios for the 2025
workforce using Equation~(\ref{eq:income-counterfactual}) and the pooled
coefficients from Equation~(\ref{eq:main}). Assigning
informal workers the formal pay prediction at their current firm size raises
their group's predicted income by 41.9 percent and total predicted income by
13.5 percent. The aggregate change is smaller because informal workers account
for 32.2 percent of baseline predicted labor income, despite representing
55.1 percent of workers. Assigning all formal workers to formal 101+ employment
raises formal workers' predicted income by 5.4 percent and total income by
3.6 percent. This is smaller than the 16.5 percent formal 101+/solo premium
in part because 54.9 percent of formal workers are already in 101+ firms and
receive no increase. The remaining workers face different estimated gaps, and
their contributions depend on their shares of baseline income.

% Generated by replication/build_counterfactual_tables.py.
\begin{table}[htbp]
  \centering
  \caption{Illustrative changes in predicted labor income, 2025}
  \label{tab:income-counterfactuals-2025}
  \small
  \begin{tabular}{>{\raggedright\arraybackslash}p{0.42\textwidth}lrr}
    \toprule
    Scenario & \shortstack[l]{Initial\\target group} & \shortstack[r]{Target group\\(\%)} & \shortstack[r]{All workers\\(\%)} \\
    \midrule
    Informal to formal, same size & Informal & 41.9 & 13.5 \\[0.5em]
    Formal workers to formal 101+ & Formal & 5.4 & 3.6 \\[0.5em]
    Informal workers to informal 101+ & Informal & 23.8 & 7.7 \\[0.5em]
    All workers to formal 101+ & All & 25.4 & 25.4 \\[0.5em]
    \bottomrule
  \end{tabular}
  \vspace{0.4em}
  \begin{minipage}{0.97\textwidth}
  \footnotesize
  Notes: Percentage changes in total predicted weekly labor income from Equation~(\ref{eq:income-counterfactual}), using the pooled coefficients from Equation~(\ref{eq:main}). The target-group denominator includes all workers initially in that group, including those already at the destination. The all-worker denominator includes the entire 2025 sample.
  \end{minipage}
\end{table}


Assigning informal workers to informal 101+ employment raises their group's
predicted income by 23.8 percent and total income by 7.7 percent. Among workers
reassigned from a smaller category to 101+, mean predicted weekly income
increases by 13.7 percent for formal workers and 24.2 percent for informal
workers. These averages exclude workers already in 101+ firms. The informal scenario
does not raise every informal worker's predicted income. The estimated informal
profile is higher in several intermediate categories than in 101+, so workers
starting in those categories receive lower predictions. Finally, assigning
everyone to formal 101+ employment raises total predicted income by
25.4 percent. This combined scenario is calculated directly. The separate
percentages cannot simply be added. Informal workers undergo both changes,
the corresponding wage ratios compound multiplicatively, and the estimated
size premium itself depends on formality.

Appendix Table~\ref{tab:income-counterfactual-sensitivity} reports alternative
destinations and an observed-income baseline for the calculation in
Equation~(\ref{eq:income-counterfactual}), retaining the coefficients from
Equation~(\ref{eq:main}). Multiplying observed income by the estimated destination-to-origin
wage ratio preserves each worker's log-wage residual. This alternative gives
similar aggregate changes. Appendix
Table~\ref{tab:income-counterfactual-support} examines how strongly the
same scenarios extrapolate beyond observed destination groups. For example,
39.1 percent of intended movers to informal 101+ employment have no observed
counterpart with the same sector, department, and occupational position in
that destination in 2025. Section~\ref{sec:discussion} discusses the
interpretation and limitations of these scenarios.

\section{Discussion}
\label{sec:discussion}

Table~\ref{tab:salaried-robustness}
shows that the larger informal 101+/solo premium in the full sample does not
generalize to salaried employment alone. This result does not imply equal
formal and informal wage levels or identical size profiles. The full sample
includes own-account work, employers, and other non-salaried employment.
Status-in-employment controls allow these groups to differ in average pay,
but Equation~(\ref{eq:main}) imposes common size coefficients across
status-in-employment categories within each formality group. Restricting the sample
to employees re-estimates all coefficients and changes the distribution of
workers across sectors, departments, and firm sizes. The resulting comparison
therefore does not identify how much of the full-sample premium is caused
by self-employment or by any other individual status-in-employment category.

The dynamic comparison in Table~\ref{tab:salaried-dynamic-comparison} also
distinguishes the two populations. The estimated informal 101+/solo premium
increases in both samples, but the employee results do not establish a
larger increase for informal than formal workers over 2008--2025. Neither employee endpoint
change over that period is significant at 5 percent under sector-department clustering.
The full-sample informal increase is significant under this procedure and
under sector clustering, though not under department or two-way clustering.

The minimum-wage distributions in
Figure~\ref{fig:minwage-distribution-formality-size} suggest one possible reason
for the smaller wage differences across intermediate formal categories.
Appendix Table~\ref{tab:minimum-wage-bunching-2025} quantifies this pattern
for 2025. Between 37.0 and 43.3 percent of formal workers in firms with
4--100 workers earn within 10 percent of the hourly benchmark. By contrast,
74.3 percent of informal solo workers and 67.0 percent of informal workers
in firms with 2--3 workers earn less than 0.9 times that benchmark.

Appendix Table~\ref{tab:minimum-wage-monthly-bunching-2025} shows that
the concentration near the minimum also appears in monthly labor
income. Between 39.3 and 47.1 percent of formal workers in firms with
4--100 workers earn within 10 percent of the monthly minimum wage. The corresponding shares are
17.9 percent for formal solo workers and 29.4 percent for formal workers in
101+ firms. The 75th percentiles for these two groups are also higher than
in the intermediate-size formal categories, at 2.11 and 2.46 monthly minimum
wages, respectively. These patterns are consistent with wage compression in
intermediate-size formal categories, but do not explain the conditional premiums
by themselves.
Neither the kernel densities nor the regressions identify the causal effect
of the minimum wage.

A wage premium that survives observable controls does not identify a single
mechanism. As \citet{BrownMedoff1989} and \citet{OiIdson1999} emphasize, the
remaining gap may reflect differences in workers' unmeasured skills or in
employers' pay practices. The latter include sharing firm profits with workers
and paying more to elicit effort or reduce turnover.
\citet{SoderbomTealWambugu2005} find substantial firm-size effects in Kenya and
Ghana after controlling for individual fixed effects. Matched employer-employee
studies separate persistent worker differences from wage differences associated
with the employer after worker fixed effects
\citep{AbowdKramarzMargolis1999,CardHeiningKline2013,
AlvarezBenguriaEngbomMoser2018,GerardLagosSeverniniCard2021,EngbomMoser2022}.
\citet{Song2019} further show that between-firm earnings inequality can grow
because high-paid workers increasingly work together, which is distinct from
an increase in what employers pay otherwise similar workers.
Recent German evidence adds a distinction between firm and establishment size.
\citet{DiegmannMullerSchoefer2026} find that the establishment-size
wage relationship disappears after adding firm fixed effects, comparing
establishments belonging to the same firm. This points to employer-wide
attributes rather than establishment size alone. GEIH does not
identify employers or follow workers across firms, so the additional controls
in our preferred specification cannot separate these explanations.

The evidence is also not a direct test of misallocation. In a
frictionless competitive benchmark with comparable labor inputs and jobs,
efficient allocation equalizes the value of the marginal product of labor
across uses. Wages then reflect that value, making persistent wage differences informative
about possible allocative frictions. If larger firms pay more for workers with
the same observed characteristics, this may indicate a higher contribution to
output or barriers to jobs in which workers could produce more. The evidence
is a marker, not proof, because
unobserved worker quality and differences in wage setting can also produce
these gaps. Structural misallocation exercises use production assumptions and
information on inputs and output to estimate marginal products and quantify
the consequences of reallocation \citep{hsieh2009misallocation}.

For Colombia, \citet{PulidoGonzalez2024}, as presented in Section 3.3 and
Annex 4 of \citet{OteroCortesEtAl2025}, use a sector-choice model in general
equilibrium, in which workers choose economic sectors and wages adjust to the
resulting labor allocation. They remove sector-specific distortions that
separate informal workers' pay from their marginal product. The model then
determines how sectoral employment shares change. The implied reallocation is
at least 35 percent of informal workers without a university degree and
41 percent of informal workers with a university degree, with an aggregate
productivity gain of 4.4 percent. These are lower bounds on movements because
offsetting worker flows need not change sectoral employment shares.
A separate exercise raises informal workers' average productivity to the formal
benchmark, yielding 15.4 percent without sectoral reallocation. Neither
exercise changes workers' formal-informal group membership, and the model does
not distinguish firm-size categories.

Our income scenarios in Section~\ref{subsec:income-counterfactuals}
assign all members of a specified group the pay
prediction for a chosen destination while keeping the estimated wage structure
fixed. Unlike the sector-choice model, they do not determine how many workers
would move or how wages and production would adjust. In particular, keeping
status in employment fixed means that a comparison for an own-account worker
does not predict earnings after becoming a salaried employee. The destination
checks in Appendix Table~\ref{tab:income-counterfactual-support} also show
that the unrestricted scenarios rely on combinations of sector, department,
status in employment, size, and formality that are sometimes sparsely observed
or absent. In the joint assignment to formal 101+ employment, 17.2 percent
of intended movers have no observed destination cell in 2025. Requiring at
least 30 observations in the pooled destination cell reduces the total income
change from 25.9 to 17.7 percent. This screen changes who is reassigned, not
the estimated income change for an identical move. The simulated income
changes are therefore not estimates of the
aggregate productivity loss from misallocation and are not directly comparable
with the output gains reported by \citet{OteroCortesEtAl2025}.

The sample comparison in Table~\ref{tab:salaried-income-comparison} shows
that the larger full-sample income changes depend substantially on including
non-salaried workers. In the joint assignment to formal 101+ employment,
restricting the income calculation to employees while retaining full-sample
coefficients reduces the predicted increase from 25.9 to 13.0 percent.
Re-estimating the coefficients for employees reduces it further to
12.3 percent. The comparison changes the population and its initial income
distribution, not just the estimated wage gaps. It also changes the ranking
of the two size-only scenarios. Moving formal workers to formal 101+
employment generates a larger total income increase than moving informal
workers to informal 101+ employment in the employee sample, whereas the
opposite holds in the full sample. The employee totals do not measure
economy-wide income changes.

Ceteris paribus, assigning all workers to formal employment in firms with
more than 100 workers raises total predicted labor income in 2025 by
25.9 percent. This is a substantial gain, but it requires an extreme
reallocation of employment. To put its magnitude in perspective, Colombia's
GDP per hour worked in 2024 was approximately 22.2 percent of the US level,
measured at purchasing power parity \citep{OECD2026Productivity}. Using
these productivity levels as a benchmark, and assuming that Colombian
output per hour increased proportionately with predicted labor income,
the ratio would rise to 27.9 percent. This would narrow the productivity
gap by 5.7 percentage points, measured relative to the US level.
Under these ceteris paribus assumptions, even a comprehensive reassignment
to formal employment in firms with more than 100 workers would leave nearly all of the international productivity
gap intact.

The policy implication is therefore not that policy should favor large firms
as such or push workers mechanically into large firms. The relevant objective
is to reduce barriers that keep productive firms small, informal, or disconnected
from larger markets, while improving workers' access to better-paid jobs.
The wage differences within informal employment support particular attention
to productive solo and micro-units. Support for their growth should be designed
as a path toward formalization, not as a subsidy for remaining informal that
weakens incentives for formal firms to remain formal. The preferred results
also support examining barriers to growth and worker movement within formal
employment, given the robust advantage of 101+ firms over intermediate formal
categories. The income scenarios also motivate examining
access to formal jobs. Holding the estimated wage structure fixed, formalization at the
current size generates a larger total income change than either of the
separate within-formality assignments to 101+ employment in both samples,
although the differences are much smaller among employees.
However, pension contribution measures worker formality, not firms' legal
registration, and the wage gaps do not establish that changing registration
alone raises productivity. These conditional rankings are not rankings of
policy returns. The scenarios do not account for implementation costs or
worker selection, nor do they establish that enough destination jobs could be created.

\section{Conclusion}

Real hourly income differs sharply
across firm-size categories, while employment remains concentrated in solo work
and micro-firms. Solo work and employment in firms with 2--10 workers together
account for about 69 percent of employment in 2008 and 61 percent in 2025.
Firms with 11--100 workers account for only 14 percent of employment in 2025.

In the preferred specification, informal workers in 101+
firms earn 37.1 percent more than observably comparable informal solo workers. Among formal workers,
the corresponding premium is 16.5 percent.
The informal contrast is statistically significant. Among informal workers, the
premium rises sharply across the smallest categories, but does not increase
monotonically through 101+. Among formal workers, the 101+ advantage over
each category from 2--3 through 51--100 is statistically significant.

In the pooled salaried-worker model, the 101+/solo premiums are similar for formal and
informal workers, at 20.2 and 20.6 percent, and a direct test does not reject
equality.
The 11--19/solo premium remains much larger among informal employees,
at 24.7 percent, than among formal employees, at 6.5 percent. The stronger
informal 101+/solo premium in the full sample therefore does not generalize
to salaried employment alone, although differences at intermediate sizes remain.

In the full sample, the estimated informal 101+/solo premium rises from 35.6 percent in 2008 to
57.7 percent in 2025, while the formal premium changes from 12.6 to
13.5 percent. The formal endpoint change is not significant.
Among employees, the informal premium rises from 20.3 to 39.1 percent and
the formal premium from 8.9 to 18.7 percent. The estimated informal widening
thus persists after restricting the sample, but its significance at 5 percent
depends on clustering in both samples. Among employees, a direct test does
not establish that the informal increase exceeds the formal increase.
These endpoint comparisons do not imply steady annual trends.

With the preferred controls, the approximate firm-size elasticity is 0.040
among formal workers and 0.097 among informal workers. Both are positive
and statistically significant, as is their difference. The formal estimate is
close in magnitude to the literature's observable-controls mean of 0.036,
while the informal estimate is about 2.7 times that benchmark.

In the 2025 income counterfactuals, at unchanged hours and an
unchanged estimated wage structure, assigning informal workers the formal
pay prediction at their current size raises total predicted weekly labor income across all workers by
13.5 percent. Assignment to 101+ employment raises total income by 3.6 percent
when restricted to formal workers and by 7.7 percent when restricted to
informal workers, holding formality fixed in both cases. Assigning both groups
to formal 101+ employment reassigns 17.02 million of the 22.6 million workers
represented by the sample, raising total predicted income by 25.4 percent and
mean predicted income among reassigned workers by 43.3 percent.

Restricting the exercise to employees and re-estimating the coefficients
reduces the joint total income increase to 11.9 percent. It reassigns
7.68 million of 12.71 million employees, with a 25.6 percent increase in
mean predicted income among those reassigned. Formalization at the current
size and assignment of informal employees to informal 101+ employment yield
smaller total changes than in the full sample, at 4.2 and 1.5 percent.
Assignment of formal employees to formal 101+ employment yields a similar
total change, at 3.5 percent. These employee totals exclude non-salaried
workers and are not economy-wide income changes.

These magnitudes summarize conditional income
differences, not productivity gains. Their interpretation is especially
limited where few or no workers share the proposed combination of sector,
department, occupational position, size, and formality.

The full-sample joint gain of 25.4 percent remains modest relative to Colombia's
productivity gap with the United States. Using the 2024 productivity levels
reported by \citet{OECD2026Productivity} as a benchmark, and assuming that
output per hour increased proportionately with predicted labor income, the
joint scenario would raise Colombia's productivity from 22.2 to 27.9 percent
of the US level, narrowing the gap by 5.7 percentage points.
Under these ceteris paribus assumptions, formalization and increased firm
scale would leave most of the international productivity gap intact.

The wage gaps motivate investigating
barriers to firm growth and access to better-paid jobs, but do not measure
the productivity gains from reallocating workers. Those gains depend on
which workers move, the jobs they can access, and the response of wages and
production to the resulting changes in employment.


\appendix
\section{Additional Tables}
\label{app:additional-tables}

This appendix reports the minimum-wage distribution tables, alternative-reference tables, comparison specifications, and income-scenario checks referenced in the main text. Table~\ref{tab:minimum-wage-bunching-2025} complements Figure~\ref{fig:minwage-distribution-formality-size} by quantifying the share of workers below, near, and above the hourly minimum-wage benchmark in 2025. Table~\ref{tab:minimum-wage-monthly-bunching-2025} reports the analogous exercise using implied monthly labor income relative to the monthly minimum wage. The alternative-reference tables use workers in firms with 101 or more workers or with 11--19 workers as the reference category. These changes alter the normalization but do not change the underlying empirical specifications. In the wage-level regressions, positive and negative size coefficients indicate higher and lower conditional wages relative to the omitted category. In the trend tables, the sign instead indicates whether that relative wage ratio rises or falls over time. In the interaction tables, the middle column is the difference between the formal and informal size coefficients, not a wage gap relative to the reference category on its own.

\begin{table}[htbp]
  \centering
  \caption{Distribution of hourly wages relative to the hourly minimum-wage benchmark, 2025}
  \label{tab:minimum-wage-bunching-2025}
  \small
  \begin{tabular}{lrrrrrrr}
    \toprule
    Firm size & Employment share & $<0.9$ & $0.9$--$1.1$ & $>1.1$ & P25 & Median & P75 \\
    \midrule
    \multicolumn{8}{l}{\textit{Formal workers}} \\
    Solo & 6.4\% & 22.2\% & 13.8\% & 64.0\% & 0.94 & 1.32 & 2.32 \\
    2--3 & 4.0\% & 21.3\% & 28.4\% & 50.3\% & 0.94 & 1.11 & 1.58 \\
    4--5 & 3.8\% & 14.2\% & 40.1\% & 45.7\% & 0.94 & 1.05 & 1.42 \\
    6--10 & 6.0\% & 11.6\% & 43.3\% & 45.1\% & 0.94 & 1.05 & 1.37 \\
    11--19 & 5.6\% & 9.0\% & 40.6\% & 50.5\% & 0.98 & 1.11 & 1.55 \\
    20--30 & 5.9\% & 7.3\% & 41.3\% & 51.4\% & 0.98 & 1.12 & 1.58 \\
    31--50 & 6.1\% & 7.5\% & 38.4\% & 54.0\% & 0.98 & 1.13 & 1.76 \\
    51--100 & 7.3\% & 6.2\% & 37.0\% & 56.8\% & 0.98 & 1.17 & 1.88 \\
    101+ & 54.9\% & 4.8\% & 28.1\% & 67.1\% & 1.02 & 1.38 & 2.53 \\
    \midrule
    \multicolumn{8}{l}{\textit{Informal workers}} \\
    Solo & 57.6\% & 74.3\% & 9.8\% & 15.9\% & 0.38 & 0.62 & 0.91 \\
    2--3 & 21.7\% & 67.0\% & 14.9\% & 18.2\% & 0.47 & 0.71 & 0.99 \\
    4--5 & 9.7\% & 61.1\% & 20.8\% & 18.1\% & 0.60 & 0.79 & 1.01 \\
    6--10 & 5.6\% & 53.6\% & 21.7\% & 24.6\% & 0.63 & 0.86 & 1.09 \\
    11--19 & 2.5\% & 49.3\% & 24.0\% & 26.6\% & 0.66 & 0.91 & 1.11 \\
    20--30 & 1.1\% & 48.3\% & 24.6\% & 27.0\% & 0.66 & 0.91 & 1.12 \\
    31--50 & 0.5\% & 47.6\% & 25.7\% & 26.7\% & 0.66 & 0.92 & 1.13 \\
    51--100 & 0.4\% & 42.3\% & 21.9\% & 35.8\% & 0.68 & 0.98 & 1.28 \\
    101+ & 0.9\% & 43.5\% & 19.5\% & 37.0\% & 0.66 & 0.95 & 1.37 \\
    \bottomrule
  \end{tabular}
  \vspace{0.3em}
  \begin{minipage}{0.95\textwidth}
  \footnotesize
  Notes: Wages are expressed relative to the hourly minimum-wage benchmark, computed as 1,423,500 pesos divided by $45\times52/12$ hours. This is an annual conversion benchmark, not a worker-specific legal compliance threshold. The columns $<0.9$, $0.9$--$1.1$, and $>1.1$ report the share within each size-formality group in each range. Employment shares are computed within each formality group. Quantiles refer to the wage-to-benchmark ratio.
  \end{minipage}
\end{table}


\begin{table}[htbp]
  \centering
  \caption{Distribution of monthly labor income relative to the monthly minimum wage, 2025}
  \label{tab:minimum-wage-monthly-bunching-2025}
  \small
  \begin{tabular}{lrrrrrrr}
    \toprule
    Firm size & Employment share & $<0.9$ & $0.9$--$1.1$ & $>1.1$ & P25 & Median & P75 \\
    \midrule
    \multicolumn{8}{l}{\textit{Formal workers}} \\
    Solo & 6.4\% & 20.1\% & 17.9\% & 62.0\% & 0.98 & 1.33 & 2.11 \\
    2--3 & 4.0\% & 9.6\% & 35.4\% & 55.0\% & 1.00 & 1.14 & 1.69 \\
    4--5 & 3.8\% & 3.4\% & 45.1\% & 51.5\% & 1.00 & 1.12 & 1.40 \\
    6--10 & 6.0\% & 2.7\% & 47.1\% & 50.2\% & 1.00 & 1.11 & 1.40 \\
    11--19 & 5.6\% & 1.8\% & 43.6\% & 54.5\% & 1.00 & 1.12 & 1.55 \\
    20--30 & 5.9\% & 1.4\% & 43.9\% & 54.7\% & 1.00 & 1.12 & 1.55 \\
    31--50 & 6.1\% & 1.4\% & 40.9\% & 57.7\% & 1.00 & 1.14 & 1.76 \\
    51--100 & 7.3\% & 1.1\% & 39.3\% & 59.6\% & 1.00 & 1.19 & 1.83 \\
    101+ & 54.9\% & 1.1\% & 29.4\% & 69.5\% & 1.00 & 1.40 & 2.46 \\
    \midrule
    \multicolumn{8}{l}{\textit{Informal workers}} \\
    Solo & 57.6\% & 78.0\% & 10.1\% & 11.9\% & 0.28 & 0.53 & 0.84 \\
    2--3 & 21.7\% & 65.5\% & 16.5\% & 18.0\% & 0.42 & 0.70 & 1.00 \\
    4--5 & 9.7\% & 56.3\% & 22.5\% & 21.3\% & 0.56 & 0.84 & 1.05 \\
    6--10 & 5.6\% & 50.9\% & 25.3\% & 23.7\% & 0.59 & 0.86 & 1.05 \\
    11--19 & 2.5\% & 46.2\% & 26.0\% & 27.7\% & 0.63 & 0.91 & 1.12 \\
    20--30 & 1.1\% & 46.3\% & 28.3\% & 25.4\% & 0.56 & 0.91 & 1.12 \\
    31--50 & 0.5\% & 42.3\% & 28.4\% & 29.3\% & 0.70 & 0.98 & 1.14 \\
    51--100 & 0.4\% & 41.9\% & 23.4\% & 34.7\% & 0.63 & 1.00 & 1.26 \\
    101+ & 0.9\% & 45.9\% & 24.9\% & 29.2\% & 0.56 & 0.95 & 1.26 \\
    \bottomrule
  \end{tabular}
  \vspace{0.3em}
  \begin{minipage}{0.95\textwidth}
  \footnotesize
  Notes: Monthly labor income refers to the main job. The columns $<0.9$, $0.9$--$1.1$, and $>1.1$ report the share within each size-formality group in each range relative to the statutory monthly minimum wage. Employment shares are computed within each formality group. Quantiles refer to the income-to-minimum ratio.
  \end{minipage}
\end{table}


\begin{table}[htbp]
  \centering
  \caption{Employment shares and mean real hourly income relative to 101+ workers, 2008 vs. 2025}
  \label{tab:descriptive-core-size-2008-2025-101-reference}
  \small
  \begin{tabular}{lrrrrrrr}
    \toprule
    & \multicolumn{2}{c}{Employment share} & \multicolumn{2}{c}{Mean hourly income} & \multicolumn{2}{c}{Wage ratio to 101+} & \multicolumn{1}{c}{Wage growth} \\
    \cmidrule(lr){2-3} \cmidrule(lr){4-5} \cmidrule(lr){6-7} \cmidrule(lr){8-8}
    Firm size & 2008 & 2025 & 2008 & 2025 & 2008 & 2025 & Annualized \\
    \midrule
    Solo & 36.1\% & 34.7\% &  5,388 &  6,534 & 0.42 & 0.39 & 1.1\% \\
    2--3 & 19.6\% & 13.8\% &  6,242 &  7,242 & 0.48 & 0.43 & 0.9\% \\
    4--5 & 7.3\% & 7.0\% &  7,196 &  8,242 & 0.56 & 0.49 & 0.8\% \\
    6--10 & 5.9\% & 5.8\% &  7,543 &  9,663 & 0.58 & 0.57 & 1.5\% \\
    11--19 & 3.4\% & 3.9\% &  8,679 & 11,634 & 0.67 & 0.69 & 1.7\% \\
    20--30 & 3.3\% & 3.3\% &  8,905 & 11,201 & 0.69 & 0.66 & 1.4\% \\
    31--50 & 2.6\% & 3.0\% & 10,424 & 12,574 & 0.81 & 0.75 & 1.1\% \\
    51--100 & 3.0\% & 3.5\% & 10,682 & 13,217 & 0.83 & 0.78 & 1.3\% \\
    101+ & 18.8\% & 25.2\% & 12,943 & 16,859 & 1.00 & 1.00 & 1.6\% \\
    \addlinespace
    All workers & 100.0\% & 100.0\% &  7,752 & 10,295 & 0.60 & 0.61 & 1.7\% \\
    \bottomrule
  \end{tabular}
  \vspace{0.3em}
  \begin{minipage}{0.95\textwidth}
  \footnotesize
  Notes: Mean hourly income is expressed in constant 2025 pesos. Wage ratios divide each firm-size mean by the mean among workers in firms with 101 or more workers in the same year. Annualized change is the compound annual change in mean hourly income between 2008 and 2025.
  \end{minipage}
\end{table}


\begin{table}[htbp]
  \centering
  \caption{Employment shares and mean real hourly income relative to 11--19 workers, 2008 vs. 2025}
  \label{tab:descriptive-core-size-2008-2025-11-19-reference}
  \small
  \begin{tabular}{lrrrrrrr}
    \toprule
    & \multicolumn{2}{c}{Employment share} & \multicolumn{2}{c}{Mean hourly income} & \multicolumn{2}{c}{Wage ratio to 11--19} & \multicolumn{1}{c}{Wage growth} \\
    \cmidrule(lr){2-3} \cmidrule(lr){4-5} \cmidrule(lr){6-7} \cmidrule(lr){8-8}
    Firm size & 2008 & 2025 & 2008 & 2025 & 2008 & 2025 & Annualized \\
    \midrule
    Solo & 36.1\% & 34.7\% &  5,388 &  6,534 & 0.62 & 0.56 & 1.1\% \\
    2--3 & 19.6\% & 13.8\% &  6,242 &  7,242 & 0.72 & 0.62 & 0.9\% \\
    4--5 & 7.3\% & 7.0\% &  7,196 &  8,242 & 0.83 & 0.71 & 0.8\% \\
    6--10 & 5.9\% & 5.8\% &  7,543 &  9,663 & 0.87 & 0.83 & 1.5\% \\
    11--19 & 3.4\% & 3.9\% &  8,679 & 11,634 & 1.00 & 1.00 & 1.7\% \\
    20--30 & 3.3\% & 3.3\% &  8,905 & 11,201 & 1.03 & 0.96 & 1.4\% \\
    31--50 & 2.6\% & 3.0\% & 10,424 & 12,574 & 1.20 & 1.08 & 1.1\% \\
    51--100 & 3.0\% & 3.5\% & 10,682 & 13,217 & 1.23 & 1.14 & 1.3\% \\
    101+ & 18.8\% & 25.2\% & 12,943 & 16,859 & 1.49 & 1.45 & 1.6\% \\
    \addlinespace
    All workers & 100.0\% & 100.0\% &  7,752 & 10,295 & 0.89 & 0.88 & 1.7\% \\
    \bottomrule
  \end{tabular}
  \vspace{0.3em}
  \begin{minipage}{0.95\textwidth}
  \footnotesize
  Notes: Mean hourly income is expressed in constant 2025 pesos. Wage ratios divide each firm-size mean by the mean among workers in firms with 11--19 workers in the same year. Annualized change is the compound annual change in mean hourly income between 2008 and 2025.
  \end{minipage}
\end{table}


\begin{table}[htbp]
  \centering
  \caption{Firm-size wage premium regressions: 101+ workers as the reference category}
  \label{tab:firm-size-wage-premium-regressions-101-reference}
  \small
  \begin{tabular}{lcccc}
    \toprule
    & \multicolumn{4}{c}{Dependent variable: log real hourly labor income} \\
    \cmidrule(lr){2-5}
    & (1) & (2) & (3) & (4) \\
    \midrule
    Firm size: Solo & -1.028*** & -0.832*** & -0.613*** & -0.372*** \\
     & (0.102) & (0.068) & (0.062) & (0.060) \\
    Firm size: 2-3 & -0.838*** & -0.584*** & -0.398*** & -0.179*** \\
     & (0.095) & (0.056) & (0.046) & (0.042) \\
    Firm size: 4-5 & -0.656*** & -0.414*** & -0.244*** & -0.056* \\
     & (0.098) & (0.052) & (0.040) & (0.032) \\
    Firm size: 6-10 & -0.522*** & -0.323*** & -0.184*** & -0.049 \\
     & (0.092) & (0.051) & (0.037) & (0.032) \\
    Firm size: 11-19 & -0.394*** & -0.247*** & -0.147*** & -0.062 \\
     & (0.081) & (0.054) & (0.041) & (0.038) \\
    Firm size: 20-30 & -0.330*** & -0.217*** & -0.139*** & -0.090** \\
     & (0.069) & (0.050) & (0.038) & (0.038) \\
    Firm size: 31-50 & -0.258*** & -0.169*** & -0.113*** & -0.090** \\
     & (0.059) & (0.045) & (0.033) & (0.034) \\
    Firm size: 51-100 & -0.199*** & -0.121*** & -0.082*** & -0.072** \\
     & (0.051) & (0.036) & (0.027) & (0.028) \\
    \midrule
    Gender, age, and education controls & No & No & Yes & Yes \\
    Formality control & No & No & No & Yes \\
    Sector-year fixed effects & No & Yes & Yes & Yes \\
    Observations & 5,426,316 & 5,426,316 & 5,426,316 & 5,426,316 \\
    $R^2$ & 0.218 & 0.277 & 0.381 & 0.391 \\
    \bottomrule
  \end{tabular}
  \vspace{0.3em}
  \begin{minipage}{0.95\textwidth}
  \footnotesize
  Notes: Specifications as in Table~\ref{tab:firm-size-wage-premium-regressions}, with 101+ workers as the omitted category. Parentheses contain sector-clustered standard errors. * $p<0.10$, ** $p<0.05$, *** $p<0.01$.
  \end{minipage}
\end{table}


\begin{table}[htbp]
  \centering
  \caption{Firm-size wage premium regressions: 11--19 workers as the reference category}
  \label{tab:firm-size-wage-premium-regressions-11-19-reference}
  \small
  \begin{tabular}{lcccc}
    \toprule
    & \multicolumn{4}{c}{Dependent variable: log real hourly labor income} \\
    \cmidrule(lr){2-5}
    & (1) & (2) & (3) & (4) \\
    \midrule
    Firm size: Solo & -0.633*** & -0.585*** & -0.466*** & -0.310*** \\
     & (0.035) & (0.041) & (0.042) & (0.051) \\
    Firm size: 2-3 & -0.444*** & -0.337*** & -0.252*** & -0.117*** \\
     & (0.029) & (0.033) & (0.028) & (0.030) \\
    Firm size: 4-5 & -0.262*** & -0.167*** & -0.097*** & 0.006 \\
     & (0.029) & (0.021) & (0.021) & (0.020) \\
    Firm size: 6-10 & -0.128*** & -0.076*** & -0.037*** & 0.013 \\
     & (0.017) & (0.010) & (0.010) & (0.010) \\
    Firm size: 20-30 & 0.064*** & 0.030*** & 0.008 & -0.028*** \\
     & (0.015) & (0.009) & (0.006) & (0.005) \\
    Firm size: 31-50 & 0.136*** & 0.078*** & 0.034*** & -0.028*** \\
     & (0.026) & (0.015) & (0.010) & (0.008) \\
    Firm size: 51-100 & 0.196*** & 0.126*** & 0.065*** & -0.010 \\
     & (0.034) & (0.023) & (0.017) & (0.014) \\
    Firm size: 101+ & 0.394*** & 0.247*** & 0.147*** & 0.062 \\
     & (0.081) & (0.054) & (0.041) & (0.038) \\
    \midrule
    Gender, age, and education controls & No & No & Yes & Yes \\
    Formality control & No & No & No & Yes \\
    Sector-year fixed effects & No & Yes & Yes & Yes \\
    Observations & 5,426,316 & 5,426,316 & 5,426,316 & 5,426,316 \\
    $R^2$ & 0.218 & 0.277 & 0.381 & 0.391 \\
    \bottomrule
  \end{tabular}
  \vspace{0.3em}
  \begin{minipage}{0.95\textwidth}
  \footnotesize
  Notes: Specifications as in Table~\ref{tab:firm-size-wage-premium-regressions}, with 11--19 workers as the omitted category. Parentheses contain sector-clustered standard errors. * $p<0.10$, ** $p<0.05$, *** $p<0.01$.
  \end{minipage}
\end{table}


\begin{table}[htbp]
  \centering
  \caption{Testing whether wage gaps relative to 101+ workers changed over time}
  \label{tab:firm-size-premium-trend-test-101-reference}
  \small
  \begin{tabular}{lcccc}
    \toprule
    Firm size & Annual trend & S.E. & $p$-value & Implied 2008--2025 change \\
    \midrule
    Solo & -0.0045 & (0.0043) & 0.313 & -7.3\% \\
    2-3 & -0.0012 & (0.0023) & 0.612 & -2.0\% \\
    4-5 & -0.0013 & (0.0023) & 0.587 & -2.1\% \\
    6-10 & -0.0023 & (0.0017) & 0.190 & -3.9\% \\
    11-19 & -0.0017 & (0.0012) & 0.175 & -2.9\% \\
    20-30 & -0.0026 & (0.0009) & 0.012 & -4.3\% \\
    31-50 & -0.0017 & (0.0009) & 0.065 & -2.9\% \\
    51-100 & -0.0018 & (0.0013) & 0.181 & -3.1\% \\
    \bottomrule
  \end{tabular}
  \vspace{0.3em}
  \begin{minipage}{0.95\textwidth}
  \footnotesize
  Notes: Annual changes in log-wage gaps from Equation~(\ref{eq:trend-benchmark}), expressed relative to 101+ workers. The final column is the compounded 2008--2025 percentage change in that wage ratio, not a percentage-point change in the premium. Two-sided tests of zero trends use sector-clustered standard errors.
  \end{minipage}
\end{table}


\begin{table}[htbp]
  \centering
  \caption{Testing whether wage gaps relative to 11--19 workers changed over time}
  \label{tab:firm-size-premium-trend-test-11-19-reference}
  \small
  \begin{tabular}{lcccc}
    \toprule
    Firm size & Annual trend & S.E. & $p$-value & Implied 2008--2025 change \\
    \midrule
    Solo & -0.0027 & (0.0039) & 0.491 & -4.5\% \\
    2-3 & 0.0005 & (0.0016) & 0.746 & 0.9\% \\
    4-5 & 0.0005 & (0.0014) & 0.736 & 0.8\% \\
    6-10 & -0.0006 & (0.0011) & 0.611 & -1.0\% \\
    20-30 & -0.0009 & (0.0010) & 0.392 & -1.5\% \\
    31-50 & 0.0000 & (0.0012) & 0.988 & 0.0\% \\
    51-100 & -0.0001 & (0.0014) & 0.940 & -0.2\% \\
    101+ & 0.0017 & (0.0012) & 0.175 & 3.0\% \\
    \bottomrule
  \end{tabular}
  \vspace{0.3em}
  \begin{minipage}{0.95\textwidth}
  \footnotesize
  Notes: Annual changes in log-wage gaps from Equation~(\ref{eq:trend-benchmark}), expressed relative to 11--19 workers. The final column is the compounded 2008--2025 percentage change in that wage ratio, not a percentage-point change in the premium. Two-sided tests of zero trends use sector-clustered standard errors.
  \end{minipage}
\end{table}


\begin{table}[htbp]
  \centering
  \caption{Firm-size coefficients by formality status: 101+ workers as reference}
  \label{tab:formality-interaction-coefficients-101-reference}
  \small
  \begin{tabular}{lccc}
    \toprule
    Firm size & $\hat{\beta}_g$ & $\hat{\theta}_g$ & $\hat{\beta}_g+\hat{\theta}_g$ \\
    \midrule
    Solo & -0.336*** & 0.182** & -0.154 \\
     & (0.046) & (0.083) & (0.088) \\
    2-3 & -0.139*** & 0.007 & -0.132*** \\
     & (0.034) & (0.047) & (0.040) \\
    4-5 & 0.007 & -0.136** & -0.129*** \\
     & (0.031) & (0.048) & (0.033) \\
    6-10 & 0.056* & -0.187*** & -0.131*** \\
     & (0.029) & (0.046) & (0.035) \\
    11-19 & 0.072** & -0.191*** & -0.119** \\
     & (0.027) & (0.044) & (0.041) \\
    20-30 & 0.066** & -0.196*** & -0.130*** \\
     & (0.024) & (0.042) & (0.040) \\
    31-50 & 0.064** & -0.175*** & -0.112*** \\
     & (0.022) & (0.041) & (0.036) \\
    51-100 & 0.073*** & -0.158*** & -0.085** \\
     & (0.021) & (0.039) & (0.030) \\
    \bottomrule
  \end{tabular}
  \vspace{0.3em}
  \begin{minipage}{0.92\textwidth}
  \footnotesize
  Notes: Log-point coefficients from Equation~(\ref{eq:formality-interaction}), with 101+ workers as the reference within each formality group and tildes suppressed. The first and third columns are informal and formal size gaps. The middle column is the formal-minus-informal difference in those size gaps. Parentheses contain sector-clustered standard errors. * $p<0.10$, ** $p<0.05$, *** $p<0.01$.
  \end{minipage}
\end{table}


\begin{table}[htbp]
  \centering
  \caption{Firm-size coefficients by formality status: 11--19 workers as reference}
  \label{tab:formality-interaction-coefficients-11-19-reference}
  \small
  \begin{tabular}{lccc}
    \toprule
    Firm size & $\hat{\beta}_g$ & $\hat{\theta}_g$ & $\hat{\beta}_g+\hat{\theta}_g$ \\
    \midrule
    Solo & -0.415*** & 0.372*** & -0.043 \\
     & (0.038) & (0.075) & (0.079) \\
    2-3 & -0.212*** & 0.191*** & -0.021 \\
     & (0.025) & (0.022) & (0.024) \\
    4-5 & -0.063*** & 0.049** & -0.013 \\
     & (0.019) & (0.018) & (0.018) \\
    6-10 & -0.016 & 0.002 & -0.014 \\
     & (0.010) & (0.012) & (0.012) \\
    20-30 & -0.005 & -0.004 & -0.009* \\
     & (0.007) & (0.009) & (0.005) \\
    31-50 & -0.009 & 0.019 & 0.010 \\
     & (0.009) & (0.012) & (0.008) \\
    51-100 & 0.004 & 0.032 & 0.036** \\
     & (0.023) & (0.026) & (0.015) \\
    101+ & -0.074** & 0.198*** & 0.124*** \\
     & (0.027) & (0.044) & (0.040) \\
    \bottomrule
  \end{tabular}
  \vspace{0.3em}
  \begin{minipage}{0.92\textwidth}
  \footnotesize
  Notes: Log-point coefficients from Equation~(\ref{eq:formality-interaction}), with 11--19 workers as the reference within each formality group and tildes suppressed. The first and third columns are informal and formal size gaps. The middle column is the formal-minus-informal difference in those size gaps. Parentheses contain sector-clustered standard errors. * $p<0.10$, ** $p<0.05$, *** $p<0.01$.
  \end{minipage}
\end{table}


Table~\ref{tab:firm-size-local-geography-robustness} adds geographic and occupational-position controls to the common-premium comparison model. Its final column absorbs sector-department-year fixed effects and controls for occupational position. It constrains the size premium to be the same for formal and informal workers. Under this restriction, workers in firms with 101 or more workers earn about 40.2 percent more than observably comparable solo workers, a statistically significant difference. Tables~\ref{tab:firm-size-local-geography-robustness-101-reference} and \ref{tab:firm-size-local-geography-robustness-11-19-reference} report the same specifications with alternative reference categories. The preferred model relaxes the common-premium restriction and is reported in Table~\ref{tab:firm-size-local-geography-formality-robustness} in the main text.

\begin{table}[htbp]
  \centering
  \caption{Common-premium estimates with geographic and occupational controls}
  \label{tab:firm-size-local-geography-robustness}
  \small
  \begin{tabular}{lcccc}
    \toprule
    & \multicolumn{4}{c}{Dependent variable: log real hourly labor income} \\
    \cmidrule(lr){2-5}
    & (1) & (2) & (3) & (4) \\
    \midrule
    Firm size: 2-3 & 0.188*** & 0.168*** & 0.167*** & 0.102*** \\
     & (0.011) & (0.011) & (0.011) & (0.013) \\
    Firm size: 4-5 & 0.307*** & 0.277*** & 0.273*** & 0.225*** \\
     & (0.014) & (0.014) & (0.014) & (0.015) \\
    Firm size: 6-10 & 0.316*** & 0.282*** & 0.278*** & 0.265*** \\
     & (0.018) & (0.017) & (0.017) & (0.016) \\
    Firm size: 11-19 & 0.302*** & 0.272*** & 0.267*** & 0.272*** \\
     & (0.018) & (0.017) & (0.017) & (0.015) \\
    Firm size: 20-30 & 0.271*** & 0.241*** & 0.237*** & 0.258*** \\
     & (0.019) & (0.017) & (0.017) & (0.015) \\
    Firm size: 31-50 & 0.271*** & 0.245*** & 0.240*** & 0.263*** \\
     & (0.018) & (0.017) & (0.017) & (0.015) \\
    Firm size: 51-100 & 0.289*** & 0.262*** & 0.257*** & 0.282*** \\
     & (0.018) & (0.017) & (0.017) & (0.015) \\
    Firm size: 101+ & 0.359*** & 0.338*** & 0.325*** & 0.338*** \\
     & (0.023) & (0.024) & (0.024) & (0.022) \\
    \midrule
    Worker and formality controls & Yes & Yes & Yes & Yes \\
    Occupational-position controls & No & No & No & Yes \\
    Fixed effects & $s\times t$ & $s\times t$, $d\times t$ & $s\times d\times t$ & $s\times d\times t$ \\
    Observations & 5,268,514 & 5,268,514 & 5,268,375 & 5,268,375 \\
    $R^2$ & 0.386 & 0.404 & 0.414 & 0.429 \\
    \bottomrule
  \end{tabular}
  \vspace{0.3em}
  \begin{minipage}{0.95\textwidth}
  \footnotesize
  Notes: The omitted category is solo workers. Column (4) corresponds to Equation~(\ref{eq:common-premium-geography}). $s$, $d$, and $t$ denote sector, department, and year. Parentheses contain sector-department-clustered standard errors. * $p<0.10$, ** $p<0.05$, *** $p<0.01$.
  \end{minipage}
\end{table}


\begin{table}[htbp]
  \centering
  \caption{Common-premium estimates with geographic and status-in-employment controls: 101+ workers as reference}
  \label{tab:firm-size-local-geography-robustness-101-reference}
  \small
  \begin{tabular}{lcccc}
    \toprule
    & \multicolumn{4}{c}{Dependent variable: log real hourly labor income} \\
    \cmidrule(lr){2-5}
    & (1) & (2) & (3) & (4) \\
    \midrule
    Firm size: Solo & -0.372*** & -0.352*** & -0.337*** & -0.343*** \\
     & (0.024) & (0.024) & (0.024) & (0.022) \\
    Firm size: 2-3 & -0.179*** & -0.178*** & -0.165*** & -0.242*** \\
     & (0.018) & (0.018) & (0.017) & (0.016) \\
    Firm size: 4-5 & -0.056*** & -0.066*** & -0.056*** & -0.117*** \\
     & (0.013) & (0.014) & (0.014) & (0.014) \\
    Firm size: 6-10 & -0.049*** & -0.063*** & -0.053*** & -0.078*** \\
     & (0.012) & (0.013) & (0.013) & (0.012) \\
    Firm size: 11-19 & -0.062*** & -0.072*** & -0.064*** & -0.071*** \\
     & (0.014) & (0.015) & (0.015) & (0.013) \\
    Firm size: 20-30 & -0.090*** & -0.100*** & -0.091*** & -0.082*** \\
     & (0.014) & (0.015) & (0.014) & (0.012) \\
    Firm size: 31-50 & -0.090*** & -0.096*** & -0.088*** & -0.076*** \\
     & (0.012) & (0.014) & (0.013) & (0.011) \\
    Firm size: 51-100 & -0.072*** & -0.079*** & -0.070*** & -0.058*** \\
     & (0.012) & (0.013) & (0.012) & (0.011) \\
    \midrule
    Worker and formality controls & Yes & Yes & Yes & Yes \\
    Status-in-employment controls & No & No & No & Yes \\
    Fixed effects & $s\times t$ & $s\times t$, $d\times t$ & $s\times d\times t$ & $s\times d\times t$ \\
    Observations & 5,426,316 & 5,426,316 & 5,426,174 & 5,426,174 \\
    $R^2$ & 0.391 & 0.410 & 0.420 & 0.435 \\
    \bottomrule
  \end{tabular}
  \vspace{0.3em}
  \begin{minipage}{0.95\textwidth}
  \footnotesize
  Notes: Specifications as in Table~\ref{tab:firm-size-local-geography-robustness}, with 101+ workers as the omitted category. $s$, $d$, and $t$ denote sector, department, and year. Parentheses contain sector-department-clustered standard errors. * $p<0.10$, ** $p<0.05$, *** $p<0.01$.
  \end{minipage}
\end{table}


\begin{table}[htbp]
  \centering
  \caption{Common-premium estimates with geographic and occupational controls: 11--19 workers as reference}
  \label{tab:firm-size-local-geography-robustness-11-19-reference}
  \small
  \begin{tabular}{lcccc}
    \toprule
    & \multicolumn{4}{c}{Dependent variable: log real hourly labor income} \\
    \cmidrule(lr){2-5}
    & (1) & (2) & (3) & (4) \\
    \midrule
    Firm size: Solo & -0.302*** & -0.272*** & -0.267*** & -0.272*** \\
     & (0.018) & (0.017) & (0.017) & (0.015) \\
    Firm size: 2-3 & -0.114*** & -0.104*** & -0.100*** & -0.169*** \\
     & (0.013) & (0.011) & (0.011) & (0.009) \\
    Firm size: 4-5 & 0.005 & 0.005 & 0.006 & -0.047*** \\
     & (0.008) & (0.008) & (0.008) & (0.007) \\
    Firm size: 6-10 & 0.014*** & 0.010** & 0.011** & -0.007 \\
     & (0.005) & (0.005) & (0.005) & (0.004) \\
    Firm size: 20-30 & -0.031*** & -0.031*** & -0.030*** & -0.014*** \\
     & (0.004) & (0.004) & (0.004) & (0.004) \\
    Firm size: 31-50 & -0.031*** & -0.027*** & -0.027*** & -0.008* \\
     & (0.005) & (0.005) & (0.005) & (0.004) \\
    Firm size: 51-100 & -0.014* & -0.010 & -0.010 & 0.010 \\
     & (0.007) & (0.007) & (0.007) & (0.007) \\
    Firm size: 101+ & 0.057*** & 0.066*** & 0.058*** & 0.066*** \\
     & (0.014) & (0.015) & (0.015) & (0.013) \\
    \midrule
    Worker and formality controls & Yes & Yes & Yes & Yes \\
    Occupational-position controls & No & No & No & Yes \\
    Fixed effects & $s\times t$ & $s\times t$, $d\times t$ & $s\times d\times t$ & $s\times d\times t$ \\
    Observations & 5,268,514 & 5,268,514 & 5,268,375 & 5,268,375 \\
    $R^2$ & 0.386 & 0.404 & 0.414 & 0.429 \\
    \bottomrule
  \end{tabular}
  \vspace{0.3em}
  \begin{minipage}{0.95\textwidth}
  \footnotesize
  Notes: Specifications as in Table~\ref{tab:firm-size-local-geography-robustness}, with 11--19 workers as the omitted category. $s$, $d$, and $t$ denote sector, department, and year. Parentheses contain sector-department-clustered standard errors. * $p<0.10$, ** $p<0.05$, *** $p<0.01$.
  \end{minipage}
\end{table}


\FloatBarrier
Table~\ref{tab:income-counterfactual-sensitivity} extends the income scenarios
in Section~\ref{subsec:income-counterfactuals}. Preserving observed log-wage
residuals gives aggregate changes close to those based on fitted income.
For example, the combined formal-101+ scenario yields 25.0 percent rather than
25.4 percent. The last two rows select the highest estimated size component
within each formality group. The selected destinations are 101+ for formal
workers and 51--100 for informal workers. Assigning informal workers to
51--100 raises total predicted income by 10.0 percent. Assigning both groups
to their respective selected destinations raises it by 13.7 percent without
changing formality. This exercise includes movement from informal 101+ to
51--100 and therefore is not uniformly a move toward larger firms. The
destinations maximize the model's fitted income within each formality group.
Because the size coefficients do not interact with other worker characteristics,
the selected size is the same for all workers of a given formality status,
not an individually estimated best job.

% Generated by replication/build_counterfactual_tables.py.
\begin{table}[p]
  \centering
  \caption{Income scenarios with alternative baselines and destinations}
  \label{tab:income-counterfactual-sensitivity}
  \small
  \begin{tabular}{>{\raggedright\arraybackslash}p{0.52\textwidth}cc}
    \toprule
    Scenario & Predicted baseline & Observed baseline \\
    & (\%) & (\%) \\
    \midrule
    Informal to formal, same size & 13.5 & 13.3 \\[0.4em]
    Formal workers to formal 101+ & 3.8 & 3.8 \\[0.4em]
    Informal workers to informal 101+ & 7.6 & 7.5 \\[0.4em]
    All workers to formal 101+ & 25.9 & 25.5 \\[0.4em]
    \midrule
    \multicolumn{3}{l}{\textit{Destinations selected from the estimated profiles}} \\
    Informal workers to informal 51--100 & 10.1 & 9.9 \\[0.4em]
    Formal to 101+, informal to 51--100 & 13.8 & 13.7 \\[0.4em]
    \bottomrule
  \end{tabular}
  \vspace{0.4em}
  \begin{minipage}{0.97\textwidth}
  \footnotesize
  Notes: Percentage changes in total weekly labor income for the entire 2025 sample. The predicted baseline uses Equation~(\ref{eq:income-counterfactual}). The observed baseline rescales observed income by the same estimated wage ratios, preserving log residuals. Both use pooled coefficients from Equation~(\ref{eq:main}). In the last two rows, destinations maximize the estimated size component within each formality group, including moves from informal 101+ to informal 51--100.
  \end{minipage}
\end{table}


Table~\ref{tab:income-counterfactual-support} reports the share of intended
movers with no observed destination counterpart and scenarios restricted by
the number of destination observations. For each intended move, we
check whether workers with the same sector, department, and occupational
position are observed in the destination size-formality category. Panel A
shows that absent destinations are especially common for assignments to
large informal categories. Panel B leaves workers at their origin whenever
the destination has fewer than the stated number of observations. Requiring
30 observations in the pooled sample reduces the total predicted income
change for assignment to informal 101+ from 7.7 to 3.7 percent. Requiring
30 observations in 2025 permits no such moves. There are 1,843 informal 101+
observations in 2025, but no destination sector-department-occupation cell
contains more than 25. The restricted scenarios thus change who is assigned
to a destination, not the estimated coefficients. Their smaller changes do
not establish smaller returns for an otherwise identical reallocation.
The screens are diagnostics of extrapolation, not tests that a job transition
is feasible or that workers are comparable on all characteristics.

% Generated by replication/build_counterfactual_tables.py.
\begin{table}[p]
  \centering
  \caption{Observed destination groups and sensitivity to extrapolation, 2025}
  \label{tab:income-counterfactual-support}
  \small
  \begin{tabular}{>{\raggedright\arraybackslash}p{0.40\textwidth}cccc}
    \toprule
    \multicolumn{5}{l}{\textit{A. Intended moves and absent destination cells (\%)}} \\
    \midrule
    Scenario & \shortstack{Intended\\movers} & \multicolumn{2}{c}{\shortstack{Absent destination\\in 2025}} & \shortstack{Absent\\in pooled} \\
    \midrule
    Informal to formal, same size & 55.1 & \multicolumn{2}{c}{6.0} & 0.4 \\[0.3em]
    Formal to 101+ & 20.2 & \multicolumn{2}{c}{6.9} & 1.8 \\[0.3em]
    Informal to 101+ & 54.6 & \multicolumn{2}{c}{39.0} & 6.9 \\[0.3em]
    All to formal 101+ & 75.3 & \multicolumn{2}{c}{17.2} & 4.6 \\[0.3em]
    Informal to 51--100 & 54.9 & \multicolumn{2}{c}{55.0} & 10.8 \\[0.3em]
    Formal to 101+, informal to 51--100 & 75.1 & \multicolumn{2}{c}{42.0} & 8.4 \\[0.3em]
    \midrule
    \multicolumn{5}{l}{\textit{B. Change in total predicted income with destination screens (\%)}} \\
    \midrule
    Scenario & No screen & \shortstack{Pooled\\$n\geq30$} & \shortstack{2025\\$n\geq10$} & \shortstack{2025\\$n\geq30$} \\
    \midrule
    Informal to formal, same size & 13.5 & 12.0 & 8.7 & 5.2\\[0.3em]
    Formal to 101+ & 3.8 & 3.2 & 2.9 & 2.4\\[0.3em]
    Informal to 101+ & 7.6 & 3.8 & 0.5 & \textit{None}\\[0.3em]
    All to formal 101+ & 25.9 & 17.7 & 11.4 & 6.7\\[0.3em]
    Informal to 51--100 & 10.1 & 2.3 & $<0.1$ & \textit{None}\\[0.3em]
    Formal to 101+, informal to 51--100 & 13.8 & 5.4 & 2.9 & 2.4\\[0.3em]
    \bottomrule
  \end{tabular}
  \vspace{0.4em}
  \begin{minipage}{0.97\textwidth}
  \footnotesize
  Notes: Destination cells are defined by sector, department, status in employment, size, and formality. Panel A reports intended movers as a percentage of all workers. The absent-destination columns instead use intended movers as their denominator. Pooled counts use all estimation years. Panel B permits a move only when its destination meets the stated unweighted observation count. Otherwise the worker stays at origin. The all-worker income denominator is unchanged. \textit{None} means that no intended move qualifies. For the joint selected-destination scenario, the 2025 $n\geq30$ column contains only formal moves.
  \end{minipage}
\end{table}

\FloatBarrier


\section*{Data and Code Availability}

Replication code and materials are available in the
\href{https://github.com/Kin-George/CJC-Monitor/tree/working-paper-v1-2026-06-13}{working-paper replication repository}.
Raw GEIH microdata are not redistributed and should be obtained directly from DANE.

\bibliographystyle{apalike}
\bibliography{Paper/references}

\end{document}
